\documentclass[11pt]{amsart}
\usepackage{amssymb,amsmath,amsthm}
\usepackage[margin=1.1in]{geometry}
\usepackage{enumitem}
\usepackage[hidelinks]{hyperref}

\newtheorem{theorem}{Theorem}[section]
\newtheorem{lemma}[theorem]{Lemma}
\newtheorem{proposition}[theorem]{Proposition}
\newtheorem{corollary}[theorem]{Corollary}
\theoremstyle{definition}
\newtheorem{definition}[theorem]{Definition}
\theoremstyle{remark}
\newtheorem{remark}[theorem]{Remark}

\newcommand{\R}{\mathbb{R}}
\newcommand{\Z}{\mathbb{Z}}
\newcommand{\frc}[1]{\{#1\}}
\DeclareMathOperator{\dist}{dist}

\title[Unbounded deficiency for unit squares at integer side]{Packing $k^2-c$ unit squares: \\ $s(k^2-c)=k$ for all large $k$}
\author{Sungjoon Ryu}
\date{Preprint, version 1.2, 7 October 2026}

\raggedbottom
\begin{document}

\begin{abstract}
Let $s(n)$ be the side of the smallest square into which $n$ non-overlapping unit squares can be packed.
We show that for every fixed $c\ge 0$ one has $s(k^2-c)=k$ for all sufficiently large integers $k$.
Equivalently, if $M(k)$ denotes the largest number of unit squares that can be placed in $[0,k]^2$ so that they are
pairwise disjoint as closed sets, then $k^2-M(k)\to\infty$; in fact $k^2-M(k)\ge \kappa\log k$ for every $k\ge2$,
with an absolute constant $\kappa>0$.
The proof combines a sharpened form of the fundamental lemma of Roth and Vaughan, with an explicit constant and a
further sharpening for nearly parallel pairs, with two simple observations: a closedness constraint on horizontal
lines, which is special to integer sides, and a quantization of heights in columns of nearly axis-parallel squares,
measured from the bottom and from the top side and organised over a continuum of scales; a count of the cracks that can
break the quantization controls the columns, and an angular argument, sharpened by a computer-assisted branch-and-bound
computation, bounds the overlap of the Roth--Vaughan rectangles at uncovered points. We obtain $\kappa=0.0353$, and
$k^2-M(k)\ge(0.0657-o(1))\log k$ as $k\to\infty$. Since $k^2-M(k)\ge1$ for every $k\ge2$ by an
elementary argument, the logarithmic bound has content only for $\log k>1/\kappa\approx28$, and the result should be
regarded as qualitative. With Daniel's unrefereed computer-assisted result $s(k^2-3)=k$ ($k\ge6$), $\kappa=0.0541$.
\end{abstract}

\maketitle

\section{Introduction}

For $n\ge 1$ let $s(n)$ denote the side of the smallest square that contains $n$ unit squares with pairwise
disjoint interiors (rotations allowed). Trivially $s(k^2)=k$, and $s(n)\le \lceil\sqrt n\,\rceil$.
For $n=k^2-c$ with $c$ small compared with $k$ it is natural to expect that one cannot do better than the trivial
$k\times k$ grid with $c$ squares removed. Nagamochi \cite{Nag} stated $s(k^2-1)=s(k^2-2)=k$ for all $k\ge2$;
the scoring assertion of his Lemma~1, on which the published proof rests, has recently been shown to fail
\cite{Che,Kar} (Karaku\c{s} points out that this does not disprove the rectangle bound that Nagamochi derives from
it), and replacement proofs have been given: of $s(k^2-1)=k$ in \cite{Kar}, through a weaker rectangle bound proved with
a strip measure, and of $s(k^2-2)=k$ (which implies $s(k^2-1)=k$), with a formalization in Lean, in \cite{Che}; see also
the discussion in \cite{Dan}. For $k^2-3$ the values $k=3,\dots,7$ are known
\cite{KS,Ben10,Ben16} ($k=5,6$ in the preprint \cite{Ben16}), and a computer-assisted proof of $s(k^2-3)=k$ for all $k\ge 6$, not yet refereed, has recently
been announced \cite{Dan} (see Corollary~\ref{cor:evand} and the paragraph after it); the same repository has since announced, again by a computer-assisted and unrefereed argument, that
$s(k^2-4)=k$ for all $k\ge5$ \cite[directory \texttt{s12/certificates/k2m4}]{Dan}. Friedman's survey \cite{Fri} contains the conjecture that $s(n^2-c)=n$ implies $s((n+1)^2-c)=n+1$, and Chung and
Graham \cite[concluding remarks]{CG09} wrote that it seems likely that $s(n^2-1)=n$, and that the same should hold for
$n^2-c$ unit squares when $c$ is fixed and $n$ is sufficiently large (attributing this to \cite{Fri}).
On the other hand, $s(k^2-c)<k$ as soon as $c$ is large in terms of $k$: Arslanov, Mustafin and Shangitbayev
\cite{AMS} show $s(n^2-n)<n$ for $n\ge 12$. Erd\H{o}s and Graham \cite{EG} packed a square of side $x$ with wasted
area $O(x^{7/11})$, and this exponent was lowered several times (see the introductions of \cite{McC,Bui}); the exponent $3/5$
was claimed by Chung and Graham \cite{CG}; an error in their argument was
pointed out by Arslanov and Bui \cite{AB}, and the bound $O(x^{3/5})$ was re-established, independently and by different methods,
in the preprints of Bui \cite{Bui} and of McClenagan \cite{McC} (the latter based on his 2024 thesis). Applied to $x$ slightly below $k$, this gives $s(k^2-Ck^{3/5})<k$ for large $k$.

Following the notes \cite[file \texttt{s12/search/FRIEDMAN.md}]{Dan}, define
\[
c^*(k):=\max\{c\in\Z_{\ge0}:\ s(k^2-c)=k\}.
\]
Since $s$ is non-decreasing, $s(k^2-c)=k$ holds exactly for $0\le c\le c^*(k)$, and Friedman's conjecture says
precisely that $c^*$ is non-decreasing. The results above give $3\le c^*(k)\le O(k^{3/5})$ for $k\ge3$; here the lower
bound uses the unrefereed announcement \cite{Dan} for $k\ge 8$ (and the later announcement would give $c^*(k)\ge4$ for $k\ge5$), and the upper bound uses the preprints \cite{McC,Bui}
(the refereed bound $O(x^{(3+\sqrt2)/7}\log x)$ of \cite{CG09} gives $c^*(k)=O(k^{(3+\sqrt2)/7}\log k)$).
Whether $c^*(k)\to\infty$, i.e.\ whether for every fixed $c$ one has $s(k^2-c)=k$ for
all large $k$, as expected in \cite{CG09}, is recorded as open in \cite{Dan}, where it is also observed that it does not follow from the lower bound
of Roth and Vaughan \cite{RV}: their theorem \cite[p.~170, Theorem]{RV} states that the wasted area $w(\alpha)$ in a square of side $\alpha$
satisfies $w(\alpha)\gg(\alpha\|\alpha\|)^{1/2}$ whenever $\alpha(\alpha-[\alpha])>1/6$, where $\|\alpha\|$ is the
distance from $\alpha$ to the nearest integer, and this vanishes as $\alpha\to k^-$. We are not aware of a proof in the
literature that $c^*(k)$ is unbounded.

\begin{theorem}\label{thm:main}
There is an absolute constant $\kappa>0$ such that for every integer $k\ge 2$, every family of $N$ unit
squares in $[0,k]^2$ that are pairwise disjoint as closed sets satisfies
\[
k^2-N\ \ge\ \kappa\log k .
\]
One may take $\kappa=0.0353$ ($\kappa=0.0319$ without the computer-assisted Lemma~\ref{lem:Kw9}, see Remark~\ref{rem:const};
both constants of this theorem use that lemma). Moreover, for every $k\ge2$ one has
$k^2-N\ge1.99954\,(\log k-13.06675)/30.418$, so that $k^2-N\ge(0.0657-o(1))\log k$ as $k\to\infty$. (Here and below $\log$ is the
natural logarithm.)
\end{theorem}

\begin{corollary}\label{cor:main}
$c^*(k)\ge \kappa\log k-1$ for $k\ge 2$. In particular $c^*(k)\to\infty$: for every fixed $c$,
$s(k^2-c)=k$ for all sufficiently large $k$. This confirms the expectation of Chung and Graham \cite{CG09} stated above.
\end{corollary}

\begin{corollary}[using an unrefereed computer-assisted result]\label{cor:evand}
Assume that $s(k^2-3)=k$ for every integer $k\ge6$, as announced, with a computer-assisted proof that has not been
refereed, in \cite[directory \texttt{s12/certificates/k2m3}]{Dan}. Then every family of $N$ unit squares in $[0,k]^2$ that are
pairwise disjoint as closed sets satisfies $k^2-N\ge4$ if $k\ge6$, and
\[
k^2-N\ \ge\ 0.0541\,\log k\qquad\text{for every integer }k\ge2 .
\]
\end{corollary}

\begin{proof}
Let $k\ge6$. Since $s(k^2-3)=k$ is not less than $k$, Lemma~\ref{lem:red} gives $k^2-3>M(k)$, so $k^2-N\ge k^2-M(k)\ge4$.
Let $F(k):=1.99954\,(\log k-13.06675)/30.418$ be the second bound of Theorem~\ref{thm:main}, and
$L_4:=13.06675+4\cdot30.418/1.99954=73.91674\ldots$, so that $F(k)\ge4$ exactly when $\log k\ge L_4$. For $2\le k\le5$,
Corollary~\ref{cor:W1} gives $k^2-N\ge1>0.0541\log5\ge0.0541\log k$. For $k\ge6$ with $\log k\le L_4$ we have $0.0541\log k\le0.0541\,L_4<3.999<4\le k^2-N$.
For $\log k>L_4$, Theorem~\ref{thm:main} gives $k^2-N\ge F(k)$, and $F(k)/\log k=\frac{1.99954}{30.418}\big(1-\frac{13.06675}{\log k}\big)$
increases with $\log k$ and equals $4/L_4>0.05411$ at $\log k=L_4$; hence $k^2-N\ge F(k)>0.05411\log k$.
\end{proof}

The assumption of Corollary~\ref{cor:evand} concerns all $k\ge6$ at once. In \cite{Dan} it is derived, for each such $k$, from a
measure on $[0,k]^2$ of total mass $k^2-4D$, $D=423621306389/500000000000>3/4$, under which every closed unit square in
$[0,k]^2$ has mass at least $1$; this property is reduced, for all $k\ge6$ simultaneously, to one finite statement about
$[0,7]^2$. The reduction and the axis-parallel case of the finite statement are checked in Lean~4 (theorems
\texttt{bentz\_of\_valid7}, \texttt{validAxis7} and \texttt{bentz\_of\_validTilt7} in \cite{Dan}), while the case of tilted squares is
certified by an exact computer search in rational arithmetic and is not formalized. The finite statement was certified again by an independently written exact
checker \cite{Wan}, and the Squares Project \cite[entry T-064]{Lev} replayed both checkers in full and rebuilt the Lean
reduction. There is no refereed publication of this result, and the recorded reviews of it are by AI agents; the repository
states that the work was produced by Claude (Anthropic) under human direction. Theorem~\ref{thm:main} does not use
\cite{Dan}; Corollary~\ref{cor:evand} changes only the constant for all $k$, not the asymptotic slope $0.0657$. The assumption is
needed only for $6\le k\le e^{58.71}$: for $\log k>13.06675+3\cdot30.418/1.99954=58.7042\ldots$ the second bound of
Theorem~\ref{thm:main} exceeds $3$, and $W$ is an integer. Through $F$, Corollary~\ref{cor:evand} also uses the computer-assisted
Lemma~\ref{lem:Kw9}; with $K_w=13$ (Remark~\ref{rem:const}) the same argument gives $4/86.0695\ldots>0.0464$.

Since $k^2-N\ge1$ for every $k\ge2$ (Corollary~\ref{cor:W1}), the bound $\kappa\log k$ says something new only for
$\log k>1/\kappa\approx28.3$; the second bound of the theorem exceeds $1$ only for $\log k>28.279$, so for smaller $k$ the
theorem is just Corollary~\ref{cor:W1} (see Remark~\ref{rem:k}).

\noindent\emph{Idea of the proof.}
On a horizontal line the squares cut out pairwise disjoint closed intervals of $[0,k]$, so at most $k-1$ of
them have length $\ge 1$ (Lemma~\ref{lem:G}). A square of inclination $a$ is cut in an interval of length
$\sec a\ge 1$ except near its top and bottom (``ramps'', of height $\sin a$). Hence, unless the line carries
some waste, it must either cross a square of large inclination, which by a sharpened form of the Roth--Vaughan fundamental lemma (Proposition~\ref{prop:FL}) costs
waste proportional to the inclination, or cross a nearly axis-parallel square through its ramp. In the latter case
we look down the column below that square: if the column is ``rigid'' (little vertical waste, only tiny
inclinations) and avoids the narrow ``ends'' of the squares in it, heights in it are quantized, i.e.\ the square's
corners sit within $1.265\cdot10^{-5}$ of integers, so its ramp cannot meet a line at a height whose fractional part is
farther than that from $0$ and $1$. A rigid column through an end of a square must pass through a tight crack between
two squares, and at most one crack per unit of height fits under a square; so the exceptional columns have measure about
$1/C$, where $1/(Cs)$ is the inclination threshold for columns at scale $s$. Thus many columns
must be non-rigid, which again costs waste. Since $k$ is an integer, the same holds for columns hanging from the top
side. A line at height $y$ is treated with the inclination threshold $1/(C'd(y))$, where $d(y)$ is the distance to the
nearer of the bottom and the top side, and each column is charged only at the least scale at which it fails to be rigid;
integrating over heights $y$ with $y_{\min}\le d(y)\lesssim k/2$, where $y_{\min}=236000$ is an absolute constant, with the
measure $dy/d(y)$ then gives a lower bound of order $\log k$; within this scheme the choice $C'=2\sqrt2$, $C=2$ is optimal.
Lines this close to the walls can be used because the excess over $1$ of the long chords on a line is itself charged to the
inclinations of the squares in the columns between them and the nearer side (Lemma~\ref{lem:Eline}), at a cost whose integral against $dy/d(y)$ is
bounded. Further refinements matter for the constant: on a line
with little waste the short chords cut from ramps add up to almost a full unit, for nearly parallel, nearly touching
squares the fundamental lemma holds with a much better constant (Proposition~\ref{prop:FLpar}), and only the uncovered points
of the Roth--Vaughan rectangles matter, where at most $13$ of them can overlap (Lemma~\ref{lem:Kw}), and at most $9$ by a
computer-assisted refinement (Lemma~\ref{lem:Kw9}). Only the numerical constants ($\kappa=0.0353$ and the explicit bound)
depend on the computer: with Lemma~\ref{lem:Kw} alone the proof gives $\kappa=0.0319$ (Remark~\ref{rem:const}).

\section{Closed packings and the reduction}

Throughout, a \emph{square} is a closed unit square. A \emph{closed packing} of $[0,k]^2$ is a finite family of squares
contained in $[0,k]^2$ that are pairwise disjoint as closed sets. Put $M(k)$ for the maximum size of a closed packing
of $[0,k]^2$ and, for a given closed packing, $W:=k^2-N$, the area of $[0,k]^2$ not covered by the squares.

\begin{lemma}\label{lem:red}
For integers $k\ge 2$ and $n\ge 1$: $s(n)<k$ if and only if $n\le M(k)$. Moreover $M(k)\le k^2-1$
\textup{(}Corollary~\ref{cor:W1} below\textup{)}, so $c^*(k)=k^2-M(k)-1$.
\end{lemma}

\begin{proof}
The minimum in the definition of $s(n)$ is attained by compactness. If $n$ squares with disjoint interiors lie in a
square of side $s<k$, place that square concentrically in $[0,k]^2$, let $O$ be the common centre, and replace each centre $c_i$ by
$O+\lambda(c_i-O)$ with $\lambda=k/s>1$, keeping orientations. If two of the original squares are separated by a line with normal $u$,
the separation $u\cdot(c_2-c_1)\ge h_1(u)+h_2(u)$ (where $h_i(u)$ is the half-width of the $i$-th square in direction $u$) becomes strict after
dilation, and every moved square stays in $[0,k]^2$: in coordinates centred at $O$, $|c_i\cdot e|+h_i(e)\le s/2$ for $e$ a
coordinate direction implies $\lambda|c_i\cdot e|+h_i(e)\le \lambda s/2-(\lambda-1)h_i(e)\le k/2$. So we obtain a closed packing of $[0,k]^2$ with $n$ squares.
Conversely, a closed packing has a positive minimum distance $\eta$ between its squares; contracting all centres
about the centre of $[0,k]^2$ by a factor $\mu<1$ close to $1$ keeps the squares disjoint and places them in a
concentric square of side $\mu k+(1-\mu)\sqrt 2<k$. Finally, $s(n)\le k$ always holds for $n\le k^2$, so
for integers $0\le c<k^2$ we have $s(k^2-c)=k$ iff $k^2-c>M(k)$.
\end{proof}

Thus Theorem~\ref{thm:main} implies Corollary~\ref{cor:main}. Note that $W$ is an integer.

\section{Elementary geometry}

The \emph{inclination} $a(S)\in[0,\pi/4]$ of a square $S$ is the least angle between a side of $S$ and a side of
$[0,k]^2$; the inclination $\theta(S_1,S_2)$ between two squares is the least angle between their sides
(this is the quotient metric on $\R/(\pi/2)\Z$, with values in $[0,\pi/4]$), and we write $S_0:=\partial[0,k]^2$ with
$\theta(S,S_0):=a(S)$. These are the definitions of \cite{RV}. Likewise, for a single side $\sigma$ of $[0,k]^2$ we put
$\theta(S,\sigma)=\theta(\sigma,S):=a(S)$.

\begin{lemma}\label{lem:geom}
Let $S$ be a square with inclination $a$ and lowest point at height $v$.
\begin{enumerate}[label=(\alph*)]
\item The vertical extent of $S$ is $[v,v+\cos a+\sin a]$.
\item For $y\in\mathcal M(S):=[v+\sin a,\ v+\cos a]$ (the \emph{middle band}, of length $\cos a-\sin a$) the line
$y=\mathrm{const}$ meets $S$ in a segment of length $\sec a\ge 1$; for $y$ in the \emph{ramp set}
$\mathcal R(S):=[v,v+\sin a)\cup(v+\cos a,\ v+\cos a+\sin a]$ the length is $<\sec a$. If $a=0$, then $\mathcal R(S)=\emptyset$.
\item Symmetrically, for $x$ in the \emph{vertical middle band} $\mathcal M'(S)$, an interval of length
$\cos a-\sin a$ \textup{(}defined in the proof; for $a=0$ it is the horizontal extent of $S$\textup{)}, the vertical line through $x$ meets $S$ in a segment of length $\sec a$, entering $S$ through a
side of slope $\pm\tan a$; the lower endpoint of that segment has height in $[v,v+\sin a]$. The complement of
$\mathcal M'(S)$ in the horizontal extent of $S$ (the \emph{vertical ramp set}) has total length $2\sin a$.
\end{enumerate}
\end{lemma}

\begin{proof}
For $a=0$ everything is clear. Let $0<a\le\pi/4$ and (up to reflection) let the vertices be
$P_0=(x_0,v)$, $P_1=P_0+(\cos a,\sin a)$, $P_3=P_0+(-\sin a,\cos a)$, $P_2=P_1+P_3-P_0$, at heights
$v,\ v+\sin a,\ v+\cos a,\ v+\cos a+\sin a$. For $y\in[v+\sin a,v+\cos a]$ the left and right boundary points lie on
the parallel sides $P_0P_3$ and $P_1P_2$, which are at distance $1$ and have direction $(-\sin a,\cos a)$; hence the
horizontal chord is $1/\cos a$. Outside this band one end of the chord lies on $P_0P_1$ or $P_3P_2$, and the chord is
shorter. For (c), put $\mathcal M'(S):=[x_0,\ x_0+\cos a-\sin a]$ (in the reflected position, its mirror image
$[x_0-\cos a+\sin a,\ x_0]$). On this interval the lower and upper boundary points lie on $P_0P_1$ and $P_3P_2$,
parallel sides at distance $1$ with direction $(\cos a,\sin a)$, so the vertical chord is $\sec a$; the lower point
has height $v+(x-x_0)\tan a\in[v,\ v+(\cos a-\sin a)\tan a]\subset[v,v+\sin a]$. The horizontal extent has length
$\cos a+\sin a$, so the vertical ramp set has length $2\sin a$.
\end{proof}

For $y\in[0,k]$ let $\omega(y)$ be the length of the part of the segment $[0,k]\times\{y\}$ not covered by squares
(the \emph{line waste}); by Fubini $\int_0^k\omega(y)\,dy=W$. Similarly, for a vertical segment $\sigma$ we call the
length of $\sigma$ not covered by squares its \emph{vertical waste}; integrating the vertical waste of
$\{x\}\times[0,k]$ over $x$ also gives $W$. Since a packing is finite, all functions of $y$ (or of $x$) considered
below are piecewise continuous with finitely many pieces, and all sets of $y$ (or of $x$) that we define are finite
unions of intervals; we shall not comment on measurability again.

\begin{lemma}[closedness]\label{lem:G}
Let $k\ge 2$. On every horizontal line, the squares of a closed packing of $[0,k]^2$ meet the segment
$[0,k]\times\{y\}$ in pairwise disjoint closed intervals, and at most $k-1$ of these have length $\ge1$.
\end{lemma}

\begin{proof}
Disjointness is inherited from the squares. If $k$ of the intervals had length $\ge 1$, their union would be a closed
subset of $[0,k]$ of full measure, hence all of $[0,k]$, which is then a disjoint union of $k\ge2$ nonempty closed
sets, contradicting connectedness.
\end{proof}

\begin{corollary}\label{cor:W1}
For $k\ge2$ every closed packing of $[0,k]^2$ has $W\ge1$; that is, $M(k)\le k^2-1$.
\end{corollary}

\begin{proof}
$W$ is an integer. If $W=0$, then $\omega(y)=0$ for some (indeed almost every) $y$, so the intervals cut out on that
line cover $[0,k]$ up to a null set; each has length $\le\sqrt2<k$, so there are at least two, and as in the proof
of Lemma~\ref{lem:G} their union is a closed set of full measure, hence $[0,k]$, a contradiction.
\end{proof}

For a square $S$ and $y\in\R$ let $c_S(y)$ be the length of the horizontal chord $S\cap(\R\times\{y\})$ (zero if the line
misses $S$).

\begin{lemma}[short chords]\label{lem:A}
Let $k\ge 2$, $y\in(0,k)$ and $0<\alpha\le\frac\pi4$, and suppose that every square meeting the line $y$ has inclination
$<\alpha$. Then
\[
\sum_{S:\ 0<c_S(y)<1}c_S(y)\ >\ 1-\omega(y)-(k-1)(\sec\alpha-1),
\]
the sum being over the squares $S$ of the packing. Moreover, for a square $S$ let $\Phi(S):=\{y:0<c_S(y)<1\}$. If
$a:=a(S)>0$, then $\Phi(S)\subset\mathcal R(S)$, $|\Phi(S)|=2\sin a\cos a$ and $\int_{\Phi(S)}c_S(y)\,dy=\sin a\cos a$; if $a(S)=0$, then
$\Phi(S)=\emptyset$.
\end{lemma}

\begin{proof}
By Lemma~\ref{lem:geom}(b), a chord of length $\ge1$ has length at most $\sec a(S)<\sec\alpha$, and by Lemma~\ref{lem:G} there are
at most $k-1$ such chords on the line; the chords cover a length $k-\omega(y)$ of $[0,k]$, which gives the inequality.
For the second part let $a>0$ and let $v$ be the height of the lowest point of $S$; in the coordinates of the proof of
Lemma~\ref{lem:geom}, for $v\le y\le v+\sin a$ the chord joins
the sides $P_0P_3$ and $P_0P_1$ and has length $(y-v)(\tan a+\cot a)=(y-v)/(\sin a\cos a)$, which is $<1$ exactly for
$y<v+\sin a\cos a$; symmetrically at the top, and in the middle band the chord is $\sec a\ge1$. Hence $\Phi(S)$ consists of two
intervals of length $\sin a\cos a$ in $\mathcal R(S)$, on each of which $c_S$ grows linearly from $0$ to $1$, so its integral
there is $\frac12\sin a\cos a$. If $a=0$, every nonzero chord has length $1$.
\end{proof}

Without an assumption on the inclinations, the same counting gives a bound in terms of the excess of the long chords.

\begin{lemma}[short chords and the excess of long chords]\label{lem:E}
Let $k\ge2$ and $y\in(0,k)$, and put
\[
E(y):=\sum_{S:\ c_S(y)\ge1}\big(c_S(y)-1\big),
\]
the sum being over the squares $S$ of the packing. Then
\[
\sum_{S:\ 0<c_S(y)<1}c_S(y)\ \ge\ 1-\omega(y)-E(y).
\]
\end{lemma}

\begin{proof}
The chords $S\cap(\R\times\{y\})$ lie in $[0,k]\times\{y\}$, are pairwise disjoint, and cover a length $k-\omega(y)$ of it. By
Lemma~\ref{lem:G}, the number $n_L$ of squares with $c_S(y)\ge1$ is at most $k-1$, and these chords have total length
$n_L+E(y)\le k-1+E(y)$. The remaining chords therefore have total length at least $k-\omega(y)-(k-1)-E(y)$.
\end{proof}

Under the assumption of Lemma~\ref{lem:A}, $E(y)<(k-1)(\sec\alpha-1)$ by Lemma~\ref{lem:geom}(b), so Lemma~\ref{lem:E} contains
the inequality of Lemma~\ref{lem:A}; Lemma~\ref{lem:Eline} below bounds $E(y)$ instead by the waste.

\section{A sharpened fundamental lemma and the cost of inclinations}

Throughout this section $k\ge4$, a closed packing of $[0,k]^2$ is fixed, $S_0:=\partial[0,k]^2$, and the \emph{waste} is the set of
points of $(0,k)^2$ that lie in no square; its area is $W$. (In the proof of Theorem~\ref{thm:main} this section and the next are used only for
$k\ge10^{12}$; for $2\le k<10^{12}$ the theorem follows from Corollary~\ref{cor:W1}, see Section~6.) We begin with two elementary facts about gaps between convex sets.

\begin{lemma}[convex gaps]\label{lem:cg}
Let $I$ be an interval of length $\lambda$, let $0\le\varphi\le1$, and let $g\ge0$ be a continuous, convex, piecewise linear
function on $I$ with finitely many break points, such that $|g'|\ge\varphi$ off the break points and the slope of $g$
increases by at least $2$ at every break point. Then $\int_Ig\ge\varphi(2-\varphi)\lambda^2/4$.
\end{lemma}

\begin{proof}
We may assume $\varphi>0$ and $\lambda>0$. Let $y^*$ be a minimum point of $g$ on the closure of $I$, and let $a$ and
$b=\lambda-a$ be the lengths of the parts of $I$ to the left and to the right of $y^*$. If $y^*$ is an end point of $I$, then
$g$ is monotone with $|g'|\ge\varphi$, so $g(y)\ge\varphi|y-y^*|$ and $\int_Ig\ge\varphi\lambda^2/2\ge\varphi(2-\varphi)\lambda^2/4$.
Otherwise $y^*$ is a break point, since $g'\ne0$ elsewhere; let $-p\le0\le q$ be the slopes of $g$ just to the left and just
to the right of $y^*$. Then $p,q\ge\varphi$ and $p+q\ge2$, and by convexity $g(y)\ge p(y^*-y)$ for $y<y^*$ and
$g(y)\ge q(y-y^*)$ for $y>y^*$. By the Cauchy--Schwarz inequality, $\lambda^2=(a+b)^2\le(pa^2+qb^2)(1/p+1/q)$, so
\[
\int_Ig\ \ge\ \frac{pa^2+qb^2}2\ \ge\ \frac{\lambda^2}{2(1/p+1/q)} .
\]
Let, say, $p\le q$. If $q\ge2-\varphi$, then $1/p+1/q\le1/\varphi+1/(2-\varphi)$. If $q<2-\varphi$, then $q\ge1$ (as
$p+q\ge2$) and $p\ge2-q$, so $1/p+1/q\le F(q):=1/(2-q)+1/q$; the function $F$ is convex on $(0,2)$ and symmetric about $1$,
so $F(q)\le F(2-\varphi)=1/\varphi+1/(2-\varphi)$. In both cases $1/p+1/q\le2/(\varphi(2-\varphi))$, and the lemma follows.
\end{proof}

Let $(s,t)$ be Cartesian coordinates in the plane. For a square $Z$ let $[\tau^-(Z),\tau^+(Z)]$ be its projection to the
$t$-axis, and for $\tau^-(Z)\le t\le\tau^+(Z)$ let $l_Z(t)$ and $r_Z(t)$ be the least and the largest $s$ with $(s,t)\in Z$.
The vertices of $Z$ at which $t=\tau^+(Z)$ or $t=\tau^-(Z)$ are its \emph{$t$-extreme} vertices.

\begin{lemma}[boundary functions]\label{lem:bd}
The function $l_Z$ is convex and $r_Z$ is concave. Each of them is piecewise linear with at most one break point, at which
the slope of $l_Z$ increases, and that of $r_Z$ decreases, by at least $2$. Off the break points $l_Z'(t)=\tan\psi$, where
$\psi\in(-\frac\pi2,\frac\pi2)$ is the angle between the $t$-axis and the side of $Z$ containing $(l_Z(t),t)$, and similarly
for $r_Z$. Moreover $\tau^+(Z)-\tau^-(Z)\ge1$.
\end{lemma}

\begin{proof}
If a side of $Z$ is parallel to the $t$-axis, then $l_Z$ and $r_Z$ are constant and $\tau^+(Z)-\tau^-(Z)=1$. Otherwise let
$w$ be the vertex of $Z$ with the least $s$. The two sides at $w$ have directions $(\cos\gamma,\sin\gamma)$ and
$(\sin\gamma,-\cos\gamma)$ for some $0<\gamma<\frac\pi2$. The point $(l_Z(t),t)$ lies on the first side for $t\ge t(w)$ and on
the second for $t\le t(w)$, where $ds/dt$ equals $\cot\gamma=\tan(\frac\pi2-\gamma)$ and $-\tan\gamma=\tan(-\gamma)$,
respectively; the slope increases at $t(w)$ by $\cot\gamma+\tan\gamma=1/(\sin\gamma\cos\gamma)\ge2$, and $l_Z$ is
convex. The case of $r_Z$ is symmetric, and $\tau^+(Z)-\tau^-(Z)=\cos\gamma+\sin\gamma\ge1$.
\end{proof}

For a square $S_1$ and $S_2$ either another square or $S_0$, let
$L$ be the segment from the centre $c_1$ of $S_1$ to the centre $c_2$ of $S_2$, respectively to a point $c_2$ of $S_0$ nearest to
$c_1$ (any choice will do), and let $R=R(S_1,S_2)$ be the closed rectangle of points at distance at most $\frac12$ from the line through $L$ whose
orthogonal projection lies on $L$. This is the rectangle used by Roth and Vaughan \cite[\S3, p.~172]{RV}; the following
proposition replaces their constant $c=10^{-10}$ by an explicit, much larger one. Let $u$ be the unit vector along $L$
and $n$ a unit normal, and use the coordinates $(s,t)$ of $c_1+su+tn$; thus $R=\{0\le s\le|L|,\ |t|\le\frac12\}$.

\begin{proposition}[sharpened fundamental lemma]\label{prop:FL}
Let $\theta:=\theta(S_1,S_2)$, and let $m'$ be the number of squares $Z\notin\{S_1,S_2\}$ of the packing that have a
$t$-extreme vertex $v\in R$ with $|t(v)|<\frac12$. Let $\mathcal T\subset(-\frac12,\frac12)$ be the set of the numbers $t(v)$
for these vertices, and let $\lambda_1,\dots,\lambda_q$ be the lengths of the intervals into which the points of $\mathcal T$
cut $(-\frac12,\frac12)$. Then the waste in $R$ has area at least
$\frac{\theta(2-\theta)}4\sum_r\lambda_r^2\ge\theta(2-\theta)/(4(m'+1))$.
\end{proposition}

\begin{proof}
Centres of squares lie in the convex set $[\frac12,k-\frac12]^2$. If $S_2$ is a square, $L$ lies in this set and $R$
lies within $\frac12$ of it, so $R\subset[0,k]^2$. If $S_2=S_0$, then $L$ is perpendicular to a side of $[0,k]^2$ nearest
to $c_1$, and $R$ is the rectangle of width $1$ standing on that side, centred at the foot of $L$; since
$c_1\in[\frac12,k-\frac12]^2$, again $R\subset[0,k]^2$, and $[0,k]^2$ lies in the half-plane $\{s\le|L|\}$.
Let $\Sigma:=\{|t|<\frac12\}$ be the open strip, $\lambda_t$ the line at height $t$, and $\ell_t:=\lambda_t\cap R$ for $|t|<\frac12$.
The points $c_1+tn$ with $|t|<\frac12$ are at distance less than $\frac12$ from $c_1$, hence interior points of $S_1$;
likewise the points $c_2+tn$ are interior points of $S_2$ if $S_2$ is a square.

(a) \emph{Squares in the strip.} Let $Z\notin\{S_1,S_2\}$ be a square that meets some chord $\ell_{t_0}$, $|t_0|<\frac12$. The
convex set $Z\cap\Sigma$ contains no interior point of $S_1$ or $S_2$, so it misses the segment $\Sigma\cap\{s=0\}$ and, if $S_2$
is a square, the segment $\Sigma\cap\{s=|L|\}$; it therefore lies on one side of each of them, and as it meets $\ell_{t_0}$, we get
$Z\cap\Sigma\subset R$ (in the wall case we also use $Z\subset\{s\le|L|\}$). Hence for $|t|<\frac12$ the square $Z$ meets $\ell_t$
if and only if $\tau^-(Z)\le t\le\tau^+(Z)$, in the coordinates $(s,t)$. If $|\tau^+(Z)|<\frac12$, a $t$-extreme vertex of $Z$
at which $t=\tau^+(Z)$ lies in $Z\cap\Sigma\subset R$, so $Z$ is one of the $m'$ squares of the statement; the same holds for
$\tau^-(Z)$. Since $\tau^+(Z)-\tau^-(Z)\ge1$ (Lemma~\ref{lem:bd}), at most one of $\tau^\pm(Z)$ lies in $(-\frac12,\frac12)$.

Consequently the numbers $\tau^\pm(Z)$ in $(-\frac12,\frac12)$, for all such $Z$, are exactly the points of $\mathcal T$, and
$\mathcal T$ has at most $m'$ elements. Cutting $(-\frac12,\frac12)$ at the points of $\mathcal T$ gives $q\le m'+1$ open intervals
$I_r$ of lengths $\lambda_r$, with $\sum_r\lambda_r=1$.

(b) \emph{The same squares in the same order.} Fix $r$. By (a), for every square $Z\notin\{S_1,S_2\}$ either $Z$ meets $\ell_t$
for all $t\in I_r$ or for none. Two disjoint compact convex sets that both meet all the lines $\lambda_t$, $t\in I_r$, are
strictly separated by a line that is not parallel to $u$, so they appear in the same order on all these lines. Hence for
$t\in I_r$ the chord $\ell_t$ meets the same squares in the same order, say $Z_0=S_1,Z_1,\dots,Z_M$, followed by $Z_{M+1}=S_2$
(or the wall $S_0$). For $1\le i\le M+1$ and $t\in I_r$ let
\[
g_i(t):=l_{Z_i}(t)-r_{Z_{i-1}}(t)\ \ge0
\]
(with $l_{Z_{M+1}}\equiv|L|$ in the wall case); this is the length of the uncovered part of $\ell_t$ between $Z_{i-1}$ and
$Z_i$, since by (a) the chord of each $Z_i$, $1\le i\le M$, on the line $\lambda_t$ lies in $\ell_t$. The regions swept by these
uncovered parts for different $i$ and $r$ are disjoint parts of the waste in $R$.

(c) \emph{Each gap is convex.} By Lemma~\ref{lem:bd}, $g_i$ is convex and piecewise linear on $I_r$, and at each of its break
points its slope increases by at least $2$. Off the break points, $g_i'=\tan\psi'-\tan\psi$, where $\psi,\psi'\in(-\frac\pi2,\frac\pi2)$
are the angles between $n$ and the sides of $Z_{i-1}$ and $Z_i$ (or the wall) bounding the gap, so
\[
|g_i'|\ \ge\ |\psi'-\psi|\ \ge\ \varphi_i:=\theta(Z_{i-1},Z_i),
\]
because the direction of a side of a square is congruent modulo $\pi/2$ to the orientation of the square. As
$\varphi_i\le\frac\pi4<1$, Lemma~\ref{lem:cg} gives $\int_{I_r}g_i\ge\varphi_i(2-\varphi_i)\lambda_r^2/4$.

(d) \emph{Summation.} By the triangle inequality for inclinations, $\sum_{i=1}^{M+1}\varphi_i\ge\theta$. Put
$\varphi_i':=\min(\varphi_i,\theta)$; then $\sum_i\varphi_i'\ge\theta$ (either some $\varphi_i\ge\theta$, or $\varphi_i'=\varphi_i$ for all $i$), and
since $x(2-x)$ is increasing on $[0,1]$,
\[
\sum_i\varphi_i(2-\varphi_i)\ \ge\ \sum_i\varphi_i'(2-\varphi_i')\ \ge\ (2-\theta)\sum_i\varphi_i'\ \ge\ \theta(2-\theta).
\]
Summing over $i$ and $r$, the waste in $R$ has area at least
\[
\sum_r\frac{\theta(2-\theta)\lambda_r^2}{4}\ \ge\ \frac{\theta(2-\theta)}{4q}\Big(\sum_r\lambda_r\Big)^2\ \ge\ \frac{\theta(2-\theta)}{4(m'+1)}.
\qedhere
\]
\end{proof}

For nearly parallel, nearly touching pairs the bound $\sum_r\lambda_r^2\ge1/(m'+1)$ is far from the truth: the cut points can
only come from squares wedged into the two small triangles left free by $S_1$ and $S_2$ in $R$, and few of them fit there.
We first record an elementary fact about sections of a square.

\begin{lemma}[sections near a vertex]\label{lem:cone}
Let $Z$ be a square, $e$ a unit vector, and $v$ a vertex of $Z$ at which $p\mapsto p\cdot e$ is smallest on $Z$. For
$0\le\nu\le\frac34$ the line $\{p:p\cdot e=v\cdot e+\nu\}$ meets $Z$ in a segment of length at least $\min(2\nu,1)$.
\end{lemma}

\begin{proof}
Measure heights by $p\cdot e-v\cdot e$, and let $0\le\psi\le\frac\pi2$ be the angle between one side at $v$ and the line
$\{p\cdot e=v\cdot e\}$; the neighbours of $v$ are at heights $\sin\psi$ and $\cos\psi$, and the opposite vertex at
$\sin\psi+\cos\psi\ge1$. If $\psi\in\{0,\frac\pi2\}$ every section at height $\nu\in(0,1]$ has length $1$. Otherwise the section at
height $\nu$ has length $\nu/(\sin\psi\cos\psi)\ge2\nu$ for $\nu\le\min(\sin\psi,\cos\psi)$, length $1/\max(\sin\psi,\cos\psi)\ge1$ for
$\nu$ between $\min(\sin\psi,\cos\psi)$ and $\max(\sin\psi,\cos\psi)$, and length $(\sin\psi+\cos\psi-\nu)/(\sin\psi\cos\psi)$ for larger $\nu$.
In the last case, if $\nu\le\frac34$, then $\sin\psi,\cos\psi\le\frac34$; with $p:=\sin\psi+\cos\psi\in[1,\sqrt2]$ we get
$p-\sin\psi\cos\psi=1-\frac12(p-1)^2\ge0.91>\nu$, so the length exceeds $1$.
\end{proof}

\begin{proposition}[nearly parallel pairs]\label{prop:FLpar}
Let $S_1$ be a square and $S_2$ a square or $S_0$, at distance less than $d$, with $\theta:=\theta(S_1,S_2)\le\theta_1$, where
$d,\theta_1\le10^{-5}$. Then $\sum_r\lambda_r^2\ge Q_*:=0.32$ in Proposition~\ref{prop:FL}, so the waste in $R(S_1,S_2)$ has
area at least $Q_*\theta(2-\theta)/4$.
\end{proposition}

\begin{proof}
Put $\varepsilon_1:=d+\theta_1\le2\cdot10^{-5}$ and $\varepsilon_0:=10^{-4}$. We use the notation of the proof of
Proposition~\ref{prop:FL}; recall that every square $Z$ contributing a point of $\mathcal T$ satisfies $Z\cap\Sigma\subset R$ and
$Z\cap(S_1\cup S_2)=\emptyset$. A point $\tau\in\mathcal T$ is \emph{deep} if $|\tau|\le\frac12-\varepsilon_0$. For each deep $\tau$
fix a square $Z$ with $\tau^-(Z)=\tau$ or $\tau^+(Z)=\tau$; we call $\tau$ an \emph{upper} point in the first case, when $Z$ lies in
$\{t\ge\tau\}$, with \emph{depth} $\frac12-\tau$, and a \emph{lower} point in the second case, with depth $\tau+\frac12$.
Different points come from different squares, since at most one of $\tau^\pm(Z)$ lies in $(-\frac12,\frac12)$. For $t$ in the
$t$-range of a set $X$ let $J_X(t)$ be the section of $X$ by the line $\lambda_t=\{p:t(p)=t\}$.

(a) \emph{Shallow points.} Let $\lambda_r'$ be the lengths of the intervals into which the deep points cut $(-\frac12,\frac12)$.
Adding the shallow points (those in $(\frac12-\varepsilon_0,\frac12)$ or $(-\frac12,-\frac12+\varepsilon_0)$) cuts only the top and
the bottom interval, say of lengths $m_+\le1$ and $m_-\le1$; parts of length at least $m_+-\varepsilon_0$ and $m_--\varepsilon_0$ remain
uncut, so $\sum_r\lambda_r^2\ge\sum_r\lambda_r'^2-4\varepsilon_0$ (if top and bottom interval coincide, the bound is
$(1-2\varepsilon_0)^2$).

(b) \emph{The wall case.} Let $S_2=S_0$. For each side of $[0,k]^2$ the distance from $S_1$ to that side equals the distance from
$c_1$ to it minus $\frac12(\cos\theta+\sin\theta)$, so the side on which $R$ stands is also a side nearest to $S_1$. In coordinates in
which $R=[-\frac12,\frac12]\times[0,|L|]$ stands on the $x$-axis and $t(p)=x$, $S_1$ has inclination $\theta$ and centre $(0,|L|)$, so it
contains the square $(0,|L|)+[-\frac{\sigma_\theta}2,\frac{\sigma_\theta}2]^2$ with $\sigma_\theta:=1/(\cos\theta+\sin\theta)\ge1-\theta$; its lowest point is
at height $<d$, so $|L|<d+\frac12(\cos\theta+\sin\theta)$ and $R\setminus S_1\subset\{y<|L|-\frac{\sigma_\theta}2\}\cup\{|x|>\frac{\sigma_\theta}2\}$, where
$|L|-\frac{\sigma_\theta}2<d+\theta\le\varepsilon_1$. If $\tau=\tau^-(Z)$ were deep, Lemma~\ref{lem:cone} would give, for $0<\nu\le\varepsilon_0/2$, a
section of $Z$ of length $\ge2\nu$ on the line $x=\tau+\nu$, which satisfies $|x|\le\frac12-\frac{\varepsilon_0}2<\frac{1-\theta}2\le\frac{\sigma_\theta}2$ (as $\theta<\varepsilon_0$); this section lies in $R\setminus S_1$,
hence in $\{y<\varepsilon_1\}$, which is impossible for $\nu=\varepsilon_0/2>\varepsilon_1/2$. The same holds for $\tau^+$. So there are no deep
points, and $\sum\lambda_r^2\ge(1-2\varepsilon_0)^2>Q_*$.

(c) \emph{Normal form.} Let now $S_2$ be a square. Two disjoint convex polygons are separated by a line parallel to a side of
one of them (the origin lies outside the convex polygon $S_1-S_2$, whose sides are parallel to sides of $S_1$ or $S_2$, so it
violates one of the inequalities defining that polygon). Exchanging $S_1$ and $S_2$ if necessary (this does not change $R$ or
$\sum\lambda_r^2$) and choosing coordinates, we may assume $S_1=[-\frac12,\frac12]^2$, $S_2\subset\{x\ge\frac12\}$, and $c_2=(a,b)$ with
$b\ge0$. Then $S_2$ has inclination $\theta$, its projections to the axes have length $w:=\cos\theta+\sin\theta\in[1,1+\theta]$, and
$S_2$ contains the axis-parallel square $S_2^\circ:=c_2+[-\frac{\sigma_\theta}2,\frac{\sigma_\theta}2]^2$, $\sigma_\theta:=1/w\ge1-\theta$. Since
$S_2\subset\{x\ge\frac12\}$ and some point of $S_2$ is within $d$ of $S_1$, we have $1\le a<1+d+\frac\theta2$, $0\le b<1+d+\frac\theta2$, and
$\ell:=a-\frac{\sigma_\theta}2$, the abscissa of the left side of $S_2^\circ$, satisfies $\frac12\le\ell<\frac12+d+\theta$. Put
$h:=b/a\in[0,1+\varepsilon_1)$ and $\rho_h:=\sqrt{1+h^2}$, so that $|L|=a\rho_h$, and take $n=(-b,a)/|L|$, so that
$t(x,y)=(ay-bx)/|L|$. A direct computation gives the levels of the vertices $V_1=(\frac12,\frac12)$, $V_2=(-\frac12,\frac12)$,
$V_3=(\frac12,-\frac12)$ of $S_1$ and $V_1'=(\ell,b-\frac{\sigma_\theta}2)$, $V_2'=(\ell,b+\frac{\sigma_\theta}2)$,
$V_3'=(\ell+\sigma_\theta,b-\frac{\sigma_\theta}2)$ of $S_2^\circ$:
\begin{gather*}
t(V_1)=\tfrac{1-h}{2\rho_h},\qquad t(V_2)=\tfrac{1+h}{2\rho_h},\qquad t(V_3)=-\tfrac{1+h}{2\rho_h},\\
t(V_1')=-\sigma_\theta\tfrac{1-h}{2\rho_h},\qquad t(V_2')=\sigma_\theta\tfrac{1+h}{2\rho_h},\qquad t(V_3')=-\sigma_\theta\tfrac{1+h}{2\rho_h}.
\end{gather*}
On the line $\lambda_t$, with $s$ as coordinate, the lines $x=\frac12$, $x=\ell$, $y=\frac12$, $y=b-\frac{\sigma_\theta}2$ are met at
$s=\frac{\rho_h}2+ht$, $s=\rho_h\ell+ht$, $s=(\frac{\rho_h}2-t)/h$ and $s=(\rho_h(b-\frac{\sigma_\theta}2)-t)/h$ respectively ($h>0$).

(d) \emph{The free gap.} For $|t|<\frac12$ the points $c_1+tn$ and $c_2+tn$ lie in the interiors of $S_1$ and $S_2$, and
$J_R(t)$ is the segment between them. Hence, if $Z$ contributes a point of $\mathcal T$ and $J_Z(t)\neq\emptyset$ with $|t|<\frac12$,
then $J_Z(t)\subset J_R(t)\setminus(S_1\cup S_2)$ lies between $J_{S_1}(t)$ and $J_{S_2}(t)\supset J_{S_2^\circ}(t)$, and
$|J_Z(t)|\le\Gamma(t)$, where $\Gamma(t)$ denotes the distance from the right end of $J_{S_1}(t)$ to the left end of
$J_{S_2^\circ}(t)$. The right end of $J_{S_1}(t)$ lies on $x=\frac12$ for $t(V_3)\le t\le t(V_1)$ and on $y=\frac12$ for
$t(V_1)\le t\le t(V_2)$; the left end of $J_{S_2^\circ}(t)$ lies on $y=b-\frac{\sigma_\theta}2$ for $t(V_3')\le t\le t(V_1')$ and on
$x=\ell$ for $t(V_1')\le t\le t(V_2')$. If $h\le1$, put $t_{\rm lo}:=t(V_1')$ and $t_{\rm hi}:=\min(t(V_1),t(V_2'))$; for
$t\in[t_{\rm lo},t_{\rm hi}]$ the two ends lie on $x=\frac12$ and $x=\ell$, so $\Gamma(t)=\rho_h(\ell-\frac12)$. If $h>1$, put
$t_{\rm lo}:=t(V_1)<0$ and $t_{\rm hi}:=t(V_1')$; for $t\in[t_{\rm lo},t_{\rm hi}]$ they lie on $y=\frac12$ and $y=b-\frac{\sigma_\theta}2$,
so $\Gamma(t)=\rho_h(b-\frac{\sigma_\theta}2-\frac12)/h$. In both cases $\Gamma\le\gamma_*:=1.415\,\varepsilon_1$ on $[t_{\rm lo},t_{\rm hi}]$, since
$\ell-\frac12<\varepsilon_1$, $b-\frac{\sigma_\theta}2-\frac12<\varepsilon_1$, $\rho_h\le1.415$, and $\rho_h/h\le\sqrt2$ for $h>1$.

(e) \emph{Depths.} Let $\tau=\tau^-(Z)$ be an upper point and suppose $\tau<t_{\rm hi}-1.5\gamma_*$. Put $t':=\tau+0.75\gamma_*$ if
this is $\ge t_{\rm lo}$, and $t':=t_{\rm lo}$ otherwise. Then $t'\in[t_{\rm lo},t_{\rm hi}]$, $|t'|<\frac12$, and
$0.75\gamma_*\le t'-\tau<\frac12$ (note $t_{\rm lo}\le0$ and $\tau\ge-\frac12+\varepsilon_0$), so by Lemma~\ref{lem:cone} and (d),
$\gamma_*\ge\Gamma(t')\ge|J_Z(t')|\ge\min(2(t'-\tau),1)\ge1.5\gamma_*$, a contradiction. Hence every upper
point has depth at most $D_u:=\frac12-t_{\rm hi}+1.5\gamma_*$, and symmetrically (using $t':=\tau-0.75\gamma_*$ or $t_{\rm hi}$) every
lower point has depth at most $D_l:=\frac12+t_{\rm lo}+1.5\gamma_*$; in particular $D_u+D_l=1-(t_{\rm hi}-t_{\rm lo})+3\gamma_*\le1+\beta_1$
with $\beta_1:=3\gamma_*$. Let $D(h):=\frac12-\frac{1-h}{2\rho_h}$; it is increasing, with derivative $(1+h)/(2\rho_h^3)$, and
$D(h)\ge\frac12$ for $h\ge1$. By (c), if $h\le1$ then $\frac12-t(V_1)=D(h)$, $\frac12-t(V_2')\le\frac12-\frac{1+h}{2\rho_h}
+\frac{1+h}{2\rho_h}\theta\le0.71\theta$ (as $(1+h)/\rho_h\le\sqrt2$ for $0\le h\le1$) and $\frac12+t(V_1')\le D(h)+\frac\theta2$, while if $h>1$ then $\frac12-t(V_1')$ and
$\frac12+t(V_1)$ are at most $\frac12\le D(h)$; so in all cases $D_u,D_l\le D(h)+0.71\theta+1.5\gamma_*\le D(h)+3\varepsilon_1$
(as $\theta\le\varepsilon_1$ and $1.5\gamma_*<2.13\,\varepsilon_1$).

(f) \emph{Two points on the same side.} Suppose $D(h)\ge0.19$; then $h\ge0.344$ (as $D(0.344)<0.19$), so that
$\tau_*:=t(V_2')-\frac12>0$ and $t(V_1)\le t(V_2')\le t(V_2)$. Let $Z$ give an upper point $\tau$ of depth $\zeta$. Then
\[
t(V_2')-\tau=\zeta+\tau_*\ \le\ D_u+\tau_*\ \le\ D(h)+\frac{1+h}{2\rho_h}-\frac12+3\varepsilon_1\ =\ \frac h{\rho_h}+3\varepsilon_1\ <\ \frac34,
\]
and $\tau\ge\frac12-D_u>-0.02$. Hence for $t\in[\tau,t(V_2')]$ each of $Z$, $S_1$, $S_2^\circ$ has a nonempty section on $\lambda_t$ (the
$t$-ranges of $S_1$ and $S_2^\circ$ contain $[t(V_3'),t(V_2')]\supset[-0.6,t(V_2')]$, and that of $Z$ contains $[\tau,\tau+1]$); the end points
of these sections depend continuously on $t$ and the three sets are pairwise disjoint, so their order on $\lambda_t$ does not
depend on $t$. By (d), applied at a level in $[\tau,\frac12)$, $J_Z(t)$ lies between $J_{S_1}(t)$ and $J_{S_2^\circ}(t)$ for all
$t\in[\tau,t(V_2')]$. At the level $t(V_2')$ the section of $S_2^\circ$ is the point $V_2'$, and that of $S_1$ ends on the side $y=\frac12$;
by (c) the gap between them is
\[
\Gamma_*:=\rho_h\Big(\ell-\frac12\Big)+\frac{\rho_h}{2h}\big(h-1+\sigma_\theta(1+h)\big)\ \le\ \rho_h+\rho_h\Big(\ell-\frac12\Big)\ <\ 2,
\]
and by Lemma~\ref{lem:cone} the section $J_Z(t(V_2'))$ has length at least $\min(2(\zeta+\tau_*),1)$. Now let two upper points of depths
$\zeta'\le\zeta$ come from squares $Z',Z$. Their sections at the level $t(V_2')$ are disjoint and lie in a segment of length
$\Gamma_*$; hence $\zeta'+\tau_*<\frac12$. If $\zeta+\tau_*\le\frac12$, then $2(\zeta'+\tau_*)+2(\zeta+\tau_*)\le\Gamma_*$, that is,
$\zeta'\le\chi_u-\zeta$ with $\chi_u:=\frac12\Gamma_*-2\tau_*$. If $\zeta+\tau_*>\frac12$, then $2(\zeta'+\tau_*)+1\le\Gamma_*$, so
\[
\zeta'\ \le\ \frac{\Gamma_*}2-t(V_2')\ =\ \frac{h(h-1)}{2\rho_h}+\frac{\rho_h}2\Big(\ell-\frac12\Big)-\frac{(1-\sigma_\theta)(1+h)(1-h)^2}{4h\rho_h}
\ <\ 1.1\,\varepsilon_1\ <\ \varepsilon_0,
\]
which is impossible for a deep point; so in this case there is at most one upper point. With
$\chi(h):=1+\frac{\rho_h}2-\frac{1+h}{\rho_h}$ we get $\chi_u\le\chi(h)+(1-\sigma_\theta)\frac{1+h}{\rho_h}+\frac{\rho_h}2(\ell-\frac12)\le\chi(h)+2.2\varepsilon_1$.
The lower points are treated in the same way below the strip, down to the level $-\frac12-\tau_*=t(V_3')$, where the section of
$S_2^\circ$ is the point $V_3'$ and that of $S_1$ ends on $x=\frac12$ (as $t(V_3)\le t(V_3')\le t(V_1)$), at distance
$\Gamma_l:=\rho_h(a-\frac{1-\sigma_\theta}2)\le\rho_h+\rho_h(a-1)$. For two lower points of depths $\zeta'\le\zeta$ this gives
$\zeta'\le\chi_l-\zeta$ if $\zeta+\tau_*\le\frac12$, where $\chi_l:=\frac12\Gamma_l-2\tau_*\le\chi(h)+(1-\sigma_\theta)\frac{1+h}{\rho_h}+
\frac{\rho_h}2(a-1)\le\chi(h)+2.2\varepsilon_1$, and $\zeta'\le\frac12\Gamma_l-\frac12-\tau_*\le\frac{h(h-1)}{2\rho_h}+\frac{\rho_h}2(a-1)+
(1-\sigma_\theta)\frac{1+h}{2\rho_h}<1.8\,\varepsilon_1<\varepsilon_0$ otherwise.

(g) \emph{Conclusion.} Let $\zeta_u\ge\zeta_u'$ ($\zeta_l\ge\zeta_l'$) be the largest and the second largest depth of an upper (lower)
point, with the value $0$ if there is no such point. Consider the intervals $(\frac12-\zeta_u+\beta_1,\frac12-\zeta_u')$,
$(-\frac12+\zeta_l',-\frac12+\zeta_l-\beta_1)$ and $(-\frac12+\zeta_l,\frac12-\zeta_u)$ (some may be empty). They contain no deep point: upper
points have depth $\zeta_u$, $\zeta_u'$ or less than $\zeta_u'$, and lower points lie at levels $\le-\frac12+D_l\le\frac12-D_u+\beta_1\le
\frac12-\zeta_u+\beta_1$; symmetrically for the second interval. They are pairwise disjoint, by the same inequality. Hence, with
$\mu_u:=\zeta_u-\zeta_u'$ and $\mu_l:=\zeta_l-\zeta_l'$,
\[
\sum_r\lambda_r'^2\ \ge\ (1-\zeta_u-\zeta_l)_+^2+(\mu_u-\beta_1)_+^2+(\mu_l-\beta_1)_+^2\ \ge\ (1-\zeta_u-\zeta_l)_+^2+\mu_u^2+\mu_l^2-2.04\beta_1,
\]
since $(\mu-\beta_1)_+^2\ge\mu^2-2\beta_1\mu$ and $\mu\le D_u,D_l\le0.51$. If $D(h)<0.19$, then $\zeta_u,\zeta_l<0.1906$ and
$\sum\lambda_r'^2\ge(1-0.3812)^2>0.38$, so by (a) $\sum_r\lambda_r^2\ge0.38-4\varepsilon_0>Q_*$. Let $D(h)\ge0.19$ and $\chi_*:=\max(\chi_u,\chi_l)$. By (f), on each side either $\mu=\zeta$ (at
most one point) or $\mu\ge(2\zeta-\chi_*)_+$; in both cases $\mu^2\ge\xi(\zeta):=\min\{\zeta^2,((2\zeta-\chi_*)_+)^2\}$. Hence
$\sum\lambda_r'^2\ge\Xi(\zeta_u,\zeta_l)-2.04\beta_1$ with $\Xi(z,z'):=(1-z-z')_+^2+\xi(z)+\xi(z')$. The function $\chi$ decreases on
$[0,h_c]$ and increases on $[h_c,\infty)$, where $h_c^3+3h_c-2=0$ (indeed $\chi'(h)=(h^3+3h-2)/(2\rho_h^3)$), so
$\chi_*\le\max(\chi(0.344),\chi(1.00002))+2.2\varepsilon_1<0.29295$. If $z+z'>1$, then $z,z'>0.49>\chi_*$ (as $z,z'\le\max(D_u,D_l)\le D(h)+3\varepsilon_1\le0.51$ by (e)) and $\Xi(z,z')\ge z^2+z'^2>0.48$.
Otherwise, choosing one of the two expressions for each of $\xi(z)$ and $\xi(z')$, $\Xi$ is bounded below by the least of:
$(1-z-z')^2+z^2+z'^2\ge\frac13$; $(1-z-z')^2+z^2+((2z'-\chi_*)_+)^2\ge\frac12(1-z')^2+((2z'-\chi_*)_+)^2\ge\frac19(2-\chi_*)^2$ (the
minimum is at $z'=\frac{1+4\chi_*}9$; for $z'\le\frac{\chi_*}2$ the bound is $\frac12(1-\frac{\chi_*}2)^2$, which is larger); and
$(1-z-z')^2+((2z-\chi_*)_+)^2+((2z'-\chi_*)_+)^2\ge\frac23(1-\chi_*)^2$ (a convex symmetric function, minimal at $z=z'=\frac{1+2\chi_*}6$).
For $\chi_*<0.29295$ all these exceed $0.3237$, so $\sum\lambda_r'^2\ge0.3237-2.04\beta_1>0.3235$ (as $\beta_1=3\gamma_*\le8.5\cdot10^{-5}$). By (a),
$\sum_r\lambda_r^2\ge0.3235-4\varepsilon_0>Q_*$.
\end{proof}

\begin{remark}
Roth and Vaughan \cite[\S3, pp.~172--173]{RV} assume that $S_1$ and $S_2$ are at distance at most $1$, bound the number $m$ of
squares meeting $R$ by $20$, and use a single strip of width $1/(4m+1)$ parallel to $L$ that contains no vertex in its
interior, and a single trapezoid in it; this gives the lower bound $(4m+1)^{-2}\tan(\theta/2m)$, about
$3.8\cdot10^{-6}\,\theta$ for $m=20$, and their lemma is stated with $c=10^{-10}$ \cite[p.~171, (3)]{RV}. In Proposition~\ref{prop:FL} the strip is cut
only where the set of squares met by the chords changes, the gaps are allowed to be convex rather than affine, and all
sub-bands and gaps are summed; with the count of Lemma~\ref{lem:count}(a) this gives $\theta(2-\theta)/40\ge\theta/33$ for pairs at
distance less than $10^{-4}$, and Proposition~\ref{prop:FLpar} gives $0.32\,\theta(2-\theta)/4\ge\theta/6.26$ for pairs at distance
less than $10^{-5}$ that are moreover nearly parallel. In randomized adversarial searches (with $S_2$ perturbed as in
Proposition~\ref{prop:FLpar}) we found no configuration with $\sum\lambda_r^2<\frac13$; the value $\frac13$ is attained by two
squares wedged into the two triangles at depth $\frac13$, and two deep points never seemed to share a triangle. Roth and Vaughan's standing
assumption $\alpha(\alpha-[\alpha])>1/6$ on the side of the container plays no role here. A version of their lemma for squares
joined by a path, with a constant derived from theirs, is given in \cite[Section~4]{BuiGood}.
\end{remark}

Let $\mathcal V_h$ consist of the squares of the packing together with the left and the right side of $[0,k]^2$. Two
distinct members $P,Q$ of $\mathcal V_h$ are \emph{consecutive on the line $y$} ($0<y<k$) if both meet the segment
$[0,k]\times\{y\}$ and the open segment between them on that line meets no square. Let $G_h$ be the graph on
$\mathcal V_h$ in which $P$ and $Q$ are adjacent if they are consecutive on a set of lines $y$ of positive measure.
Define $\mathcal V_v$ and $G_v$ in the same way, with vertical lines and the bottom and the top side of $[0,k]^2$.

\begin{lemma}[visibility graphs are planar]\label{lem:planar}
The graphs $G_h$ and $G_v$ are planar.
\end{lemma}

\begin{proof}
We treat $G_h$. For every edge $e=\{P,Q\}$ choose a line $y_e\in(0,k)$ on which $P$ and $Q$ are consecutive, with
$y_e\ne y_{e'}$ for $e\ne e'$ (each edge has a set of such lines of positive measure). Let $P$ be to the left of $Q$ on
$y_e$, and let $a_e\in P$ and $b_e\in Q$ be the end points of the segment between them; $a_e\ne b_e$ if both are
squares, since they are disjoint (if one of them is a side, possibly $a_e=b_e$, and the middle piece of the path below is empty). Represent a square by its centre, the left side by $o_-:=(-2,k/2)$ and the right side
by $o_+:=(k+2,k/2)$, and draw $e$ as the path from the representative of $P$ to $a_e$, along $[a_e,b_e]$, and on to the
representative of $Q$. Here the piece inside a square is the segment from its centre to the boundary point; it lies in
the square, and all its points except the last one are interior points. The piece for the left side is the segment
from $o_-$ to $(-1,y_e)$ followed by the segment from $(-1,y_e)$ to $(0,y_e)=a_e$, and symmetrically for the right side.
Two such paths meet only in common end points. Indeed, segments from the centre of a square to distinct boundary points
(the points $a_e$, $b_e$ of different edges lie on different lines) meet only at the centre, and different squares are
disjoint. The open segments $(a_e,b_e)$ lie on distinct lines and meet no square. The pieces for the sides lie outside
$(0,k)\times\mathbb R$, the horizontal ones lie on distinct lines, and the segments from $o_-$ to the distinct points
$(-1,y_e)$ meet only at $o_-$ (likewise at $o_+$). So this is a plane drawing of $G_h$.
\end{proof}

For $0<d\le10^{-4}$ put $\Lambda=\Lambda(d):=d+\sqrt2$ and $R_0=R_0(d):=\sqrt{\Lambda^2+\frac14}<1.5001$. We use the
following theorem of Fodor \cite[Theorem~1]{Fod}: \emph{the smallest circle that contains $12$ points with mutual distances at least $1$
has radius $R_{12}=1.5148\ldots$} (explicitly, $R_{12}=1/(\sqrt3\,x_0)$, where $x_0$ is the smallest positive root of
$9x^5-15x^4+7x^3-3x+1$). We also use Oler's inequality \cite[Theorem~1]{Ole} in the following Euclidean form, which is
inequality~(1) of Folkman and Graham \cite{FG}, where it also receives a short proof: \emph{a compact convex set $X\subset\R^2$
contains at most $\frac2{\sqrt3}A(X)+\frac12P(X)+1$ points with mutual distances at least $1$}, where $A$ and $P$ denote area and
perimeter. (Oler's theorem is stated for the Minkowski distance of a convex, centrally symmetric Jordan curve, with the area
divided by the critical determinant, which is $\sqrt3/2$ for the Euclidean norm, and for finite sets that lie in the closed region
bounded by a Jordan polygon all of whose vertices belong to the set. The Euclidean form above follows by applying it to the boundary
of the convex hull of the points, whose area and perimeter do not exceed those of $X$; collinear points are trivial. Folkman and
Graham \cite[p.~751]{FG} make the same convex-hull reduction from their own inequality for planar simplicial complexes, which gives an
independent proof of the form above.) Put
\[
K:=27,\qquad q_0:=10 .
\]

\begin{lemma}[counting]\label{lem:count}
Let $0<d\le10^{-4}$.
\begin{enumerate}[label=(\alph*)]
\item If $S_1$ is a square and $S_2$ is a square or $S_0$, at distance $<d$, then $m'\le q_0-1=9$ in
Proposition~\ref{prop:FL}; hence the waste in $R(S_1,S_2)$ has area at least $\theta(2-\theta)/(4q_0)$, $\theta=\theta(S_1,S_2)$.
\item Let $k\ge4$ and let $G$ be $G_h$ or $G_v$. For an edge $e=\{P,Q\}$ of $G$ whose members are at distance $<d$ put $R_e:=R(P,Q)$ if
both are squares, and $R_e:=R(P,S_0)$ if $P$ is a square and $Q$ a side. Then every point lies in $R_e$ for at most $K$
such edges $e$ \textup{(}the two sides of $[0,k]^2$ that are vertices of $G$ are at distance $k>d$\textup{)}.
\end{enumerate}
\end{lemma}

\begin{proof}
Every point of a unit square is within $\frac{\sqrt2}2$ of its centre, so the centres of two squares at distance $<d$ are
at distance $<\Lambda$, and the centre of a square at distance $<d$ from $S_0$ is at distance $<d+\frac{\sqrt2}2$ from
$S_0$.

(a) For each of the $m'$ squares $Z$ of Proposition~\ref{prop:FL} choose a vertex $v_Z$ as there, and let
$Q_Z:=Z\cap B(v_Z,1)$; in coordinates in which $Z=[0,1]^2$ and $v_Z=0$ this is the quarter-disc $\{x,y\ge0,\ x^2+y^2\le1\}$,
of area $\frac\pi4$. The sets $Q_Z$, $S_1$ and (if it is a square) $S_2$ are pairwise disjoint, and they lie in the closed
$1$-neighbourhood $R+B(1)$ of $R$, because $v_Z,c_1,c_2\in R$ and $S_i\subset B(c_i,\frac{\sqrt2}2)$. By Steiner's formula,
$R+B(1)$ has area $|L|+2(|L|+1)+\pi$. If $S_2$ is a square, then $|L|<\Lambda$ and
$\frac\pi4m'+2\le3|L|+2+\pi$, so $m'\le\frac{12}\pi|L|+4<\frac{12}\pi\Lambda+4<10$. If $S_2=S_0$, then $|L|<d+\frac{\sqrt2}2$, and
all these sets lie in $[0,k]^2$, hence in the closed half-plane $\mathcal H$ that is bounded by the line through the side of
$[0,k]^2$ on which $R$ stands and contains $[0,k]^2$. The part of $R+B(1)$ in $\mathcal H$ has area $3|L|+1+\frac\pi2$ (we remove a
$1\times1$ rectangle and two quarter-discs of radius $1$ beyond that line), whence $\frac\pi4m'+1\le3|L|+1+\frac\pi2$ and
$m'\le\frac{12}\pi|L|+2<5$. Now apply Proposition~\ref{prop:FL}.

(b) If $p\in R_e$, the projection of $p$ on the line of $L$ lies on $L$, so $p$ is within $\sqrt{|L|^2+\frac14}$ of both
end points of $L$. For an edge joining two squares, both centres therefore lie in the open disc $B(p,R_0)$. For an edge
joining a square $P$ and a side $Q$, the square is at distance $<d$ from $S_0\supset Q$, so (a) applies to $(P,S_0)$, and the
centre of $P$ lies within $\sqrt{(d+\frac{\sqrt2}2)^2+\frac14}<R_0$ of $p$. Disjoint squares contain the disjoint open discs of
radius $\frac12$ about their centres, so the centres are at mutual distance at least $1$; as $R_0<R_{12}$, at most $11$
squares have their centre in $B(p,R_0)$. Recall that a simple planar graph with $v$ vertices has at most $\max(1,3v-6)$ edges, and that the edges in question are edges of the planar graph $G$ (Lemma~\ref{lem:planar}).
If none of them contains a side of $[0,k]^2$, they span at most $11$ vertices, so there are at most $3\cdot11-6=27=K$ of them.

Otherwise let $e_0=\{P_0,Q_0\}$ be one of them with $Q_0$ a side, and let $\sigma_0$ be the side of $[0,k]^2$ on which
$R_{e_0}=R(P_0,S_0)$ stands. Then $p$ is at distance at most $|L_{e_0}|<d+\frac{\sqrt2}2$ from the line through $\sigma_0$, while every
centre is at distance at least $\frac12$ from that line, on the same side as $[0,k]^2$. Hence, if $\nu$ is the unit normal of $\sigma_0$
pointing away from $[0,k]^2$, every centre $z$ in question satisfies $(z-p)\cdot\nu<c_0:=d+\frac{\sqrt2}2-\frac12<0.20721$, so all these
centres lie in the compact convex set $D_0:=\{z\in\overline B(p,R_0):(z-p)\cdot\nu\le c_0\}$. With
$\phi:=\arccos(c_0/R_0)$, this set has area $R_0^2(\pi-\phi)+c_0\sqrt{R_0^2-c_0^2}<4.155$ and perimeter $2R_0(\pi-\phi)+2\sqrt{R_0^2-c_0^2}<8.1$,
so by Oler's inequality it contains at most $\frac2{\sqrt3}\cdot4.155+\frac{8.1}2+1<9.85$, that is, at most $9$, centres.
Moreover, only one side of $[0,k]^2$ occurs among the edges in question: if $\{P,Q\}$ is one of them with $Q$ the left side, then the
centre of $P$ is within $d+\frac{\sqrt2}2$ of the left side and $p$ is within $\sqrt{(d+\frac{\sqrt2}2)^2+\frac14}$ of that centre, so $p$ is
within $1.58$ of the left side; likewise for the right side (or for the bottom and the top side if $G=G_v$), and $k\ge4>3.16$. So the
edges span at most $9+1$ vertices, and there are at most $3\cdot10-6=24<K$ of them.
\end{proof}

Only the uncovered points of the rectangles $R_e$ will matter below, and there the multiplicity is much smaller.

\begin{lemma}[multiplicity at uncovered points]\label{lem:Kw}
Let $0<d\le10^{-4}$, let $G$ be $G_h$ or $G_v$, and define $R_e$ for the edges $e$ of $G$ whose members are at distance $<d$ as
in Lemma~\ref{lem:count}(b). Then every point of the waste lies in $R_e$ for at most $K_w:=13$ such edges.
\end{lemma}

\begin{proof}
Let $p$ be a point of the waste, $\Lambda:=d+\sqrt2$, and call an edge $e$ of $G$ \emph{relevant} if its members are at distance
$<d$ and $p\in R_e$; let $L_e$ be the segment defining $R_e$. Every square contains the closed disc of radius $\frac12$ about its
centre, so every centre is at distance $>\frac12$ from $p$. As in the proof of Lemma~\ref{lem:count}(b), for a relevant $e$ the
orthogonal projection of $p$ on the line of $L_e$ lies on $L_e$, and $p$ is at distance $\le\frac12$ from that line. If both
members of $e$ are squares, then $|L_e|<\Lambda$ and $p$ is within $\sqrt{(\Lambda/2)^2+\frac14}$ of the nearer end point of $L_e$;
if $e$ joins a square $P$ and a side, then $|L_e|<d+\frac{\sqrt2}2$ and $p$ is within $\sqrt{(d+\frac{\sqrt2}2)^2+\frac14}$ of the centre
of $P$. As $\Lambda/2\le d+\frac{\sqrt2}2$, in both cases some member of $e$ is a square whose centre lies within
$\rho_1:=\sqrt{(d+\frac{\sqrt2}2)^2+\frac14}<0.8662$ of $p$; we fix one such member and call its centre the \emph{inner centre}
of $e$. Put $f(r_1,r_2):=(r_1^2+r_2^2-1)/(2r_1r_2)$; if two points are at distances $r_1,r_2$ from a point $z$ and at distance
$\ge1$ from each other, the cosine of the angle they subtend at $z$ is at most $f(r_1,r_2)$.

(a) \emph{Degree at a centre.} Let $c$ be the centre of a square $P$, $r:=|p-c|$, and let $e_1,\dots,e_m$ be the relevant edges
containing $P$. Let $c_j$ be the end point of $L_{e_j}$ other than $c$: the centre of the other member, or, if the other member
is a side, the point of $S_0$ used to define $R_{e_j}=R(P,S_0)$. The angle $\psi_j$ between $p-c$ and $c_j-c$ satisfies
$\cos\psi_j\ge0$ and $r\sin\psi_j\le\frac12$, so $\psi_j\le\arcsin\frac1{2r}$, and all directions $c_j-c$ lie in an arc of length
$2\arcsin\frac1{2r}$. At most one $e_j$ joins $P$ to a side: the two sides of $[0,k]^2$ that are vertices of $G$ are at distance $k$,
while $P$ has width at most $\sqrt2<k-2d$ in the perpendicular direction. If $c_i$ and $c_j$ are centres, they are at distances in
$[1,\Lambda)$ from $c$ (squares at distance $<d$) and at distance $\ge1$ from each other. For fixed $r_2\ge1$ the function
$r_1\mapsto f(r_1,r_2)=\frac{r_1}{2r_2}+\frac{r_2^2-1}{2r_2}\cdot\frac1{r_1}$ is convex, and symmetrically in $r_2$, so the maximum of $f$ on
$[1,\Lambda]^2$ is attained at a corner, where $f$ takes the values $\frac12$, $\frac\Lambda2$ and $1-\frac1{2\Lambda^2}$; hence the directions of
$c_i-c$ and $c_j-c$ differ by at least $\gamma_0:=\arccos(1-\frac1{2\Lambda^2})$. If $c_i$ lies on a side $\sigma$ of $[0,k]^2$, then $c_i-c$ is
a positive multiple of the unit normal $\nu$ of $\sigma$ pointing away from $[0,k]^2$, $c$ is at distance less than $d+\frac{\sqrt2}2$
from the line of $\sigma$, and every centre $c_j$ is at distance at least $\frac12$ from it, on the same side; so
$(c_j-c)\cdot\nu<d+\frac{\sqrt2}2-\frac12$, and as $|c_j-c|\ge1$ the two directions differ by more than
$\arccos(d+\frac{\sqrt2}2-\frac12)>78^\circ>\gamma_0$. Consequently
\[
m\ \le\ D(r):=1+\Big\lfloor\frac{2\arcsin(1/(2r))}{\gamma_0}\Big\rfloor ,
\]
and $D(r)\le5$ for $r>\frac12$, since $\pi<5\gamma_0$.

(b) \emph{Counting.} Let $c_1,\dots,c_n$ be the distinct inner centres, at distances $r_i\in(\frac12,\rho_1]$ from $p$. Every relevant
edge contains the square of its inner centre, so by (a) there are at most $\sum_iD(r_i)$ relevant edges. Put $m_i:=\max(2,D(r_i))$, and
$R_2:=\rho_1$, $R_m:=\min\big(\rho_1,1/(2\sin\frac{(m-1)\gamma_0}2)\big)$ for $m=3,4,5$; then $r_i\le R_{m_i}$. For $0<r_1,r_2<1$ we have
$\partial f/\partial r_1=(r_1^2-r_2^2+1)/(2r_1^2r_2)>0$, and symmetrically, so two inner centres $c_i,c_j$ (at distance $\ge1$) are seen
from $p$ at an angle at least $\gamma(m_i,m_j)$, where $\gamma(a,b):=\arccos f(R_a,R_b)$. Order the inner centres by their direction
from $p$. If $n\ge3$, the angles between consecutive directions add up to $2\pi$, so $\sum_{i}\gamma(m_i,m_{i+1})\le2\pi$ (indices
modulo $n$); if $n\le2$, then $\sum_im_i\le10$. Since $R_e$ does not depend on $d$, every edge that is relevant for $d$ is relevant
for $10^{-4}$, so we may take $d=10^{-4}$, where
$\gamma_0>41.406^\circ$, $\rho_1<0.86611$, $R_3<0.75598$, $R_4<0.56571$, $R_5<0.50396$, and $\gamma(a,b)\ge\gamma(2,2)>70.52^\circ$, so
$n\le5$. A finite check of all sequences $(m_1,\dots,m_n)\in\{2,3,4,5\}^n$, $n\le5$, shows that every sequence with
$\sum_im_i\ge15$ violates these angle conditions, with a margin of more than $25^\circ$, and that the only sequences with
$\sum_im_i=14$ that satisfy them are, up to rotation, $(5,2,5,2)$ (which uses about $359^\circ$). The check, in $50$-digit
arithmetic, is part of the accompanying script. So there are at most $14$ relevant edges, and at most $13$ unless two inner
centres have class $5$.

(c) \emph{Two centres of class $5$.} For an inner centre $c$ write $\deg(c)$ for the number of relevant edges containing its square;
the number of relevant edges is at most $\sum_i\deg(c_i)$, and an edge both of whose members are squares with centres among
$c_1,\dots,c_n$ is counted twice in this sum. Suppose that two inner centres, say $c_1$ and $c_3$, have class $5$, i.e.\ $r_1,r_3\le R_5<0.50396$. Then
$1\le|c_1-c_3|\le2R_5<1.008$ and the angle at $p$ between $c_1-p$ and $c_3-p$ is at least $\gamma(5,5)>165.6^\circ$, so the angles of
the triangle $pc_1c_3$ at $c_1$ and at $c_3$ are acute, and by the law of sines the angle at $c_1$ between $p-c_1$ and $c_3-c_1$ is at
most $\psi:=\arcsin(R_5\sin(\pi-\gamma(5,5)))<7.19^\circ$. If $c_3$ is
the other end point of a relevant edge at $c_1$, this edge is counted in $\deg(c_1)$ and in $\deg(c_3)$, so there are at most
$\sum_iD(r_i)-1$ relevant edges. Otherwise let $c_j$ be the other end point of a relevant edge at $c_1$. If $c_j$ is a centre, it is
the centre of a square other than that of $c_3$, so $|c_j-c_3|\ge1$; as $|c_j-c_1|$ and $|c_3-c_1|$ lie in $[1,\Lambda)$ (note $2R_5<\Lambda$),
the argument of (a) shows that the directions of $c_j-c_1$ and $c_3-c_1$ differ by at least $\gamma_0$. If $c_j$ lies on a side, they
differ by more than $78^\circ>\gamma_0$, as in (a). Hence the directions of the $c_j-c_1$ lie in the arc of half-width
$\arcsin\frac1{2r_1}<90^\circ$ about the direction of $p-c_1$ from which the open arc of half-width $\gamma_0$ about the direction of
$c_3-c_1$ has been removed, that is, in two closed arcs, each of length less than $90^\circ+7.19^\circ-41.406^\circ<2\gamma_0$. Each of them
contains at most two of these directions, so $\deg(c_1)\le4=D(r_1)-1$, and again there are at most $\sum_iD(r_i)-1$ relevant
edges. In all cases, by (b), there are at most $\sum_iD(r_i)-1\le\sum_im_i-1\le13$ relevant edges if two inner centres have class
$5$, and at most $13$ otherwise.
\end{proof}

The proof of Lemma~\ref{lem:Kw} uses no property of $G$ except that the only sides involved are two opposite ones; the bound
holds for any family of such pairs at distance $<d$. A computer-assisted refinement of the same counting gives the bound $9$, which
is sharp for the number of all pairs at distance $<d$ (see the remark after the proof), though not known to be sharp for the edges
of a single graph $G$.

\begin{lemma}[multiplicity at uncovered points, computer-assisted]\label{lem:Kw9}
Under the assumptions of Lemma~\ref{lem:Kw}, every point of the waste lies in $R_e$ for at most $K_w^*:=9$ of the edges $e$
considered there.
\end{lemma}

\begin{proof}
As in the proof of Lemma~\ref{lem:Kw} we may take $d=10^{-4}$ and $p=0$, and we keep the notation $\Lambda$, $\rho_1$, $\gamma_0$.
Let $I$ be the set of \emph{all} centres within distance $\rho_1$ of $p$; they satisfy $\frac12<|c|\le\rho_1$ and are pairwise at
distance at least $1$, so $n:=|I|\le5$. By the proof of Lemma~\ref{lem:Kw}, every relevant edge contains a square with centre in
$I$. At most $E_{II}$ relevant edges join two squares with centres in $I$, where $E_{II}$ is the number of pairs in $I$ at distance
less than $\Lambda$. Every other relevant edge joins a square with centre $c_i\in I$ either to a side of $[0,k]^2$ or to a square whose
centre $x$ satisfies $|x|>\rho_1$, $1\le|x-c_i|<\Lambda$, $p\in R(c_i,x)$ and $|x-c_j|\ge1$ for all $c_j\in I\setminus\{c_i\}$, where
$R(c_i,x)$ denotes the rectangle $R$ of the two squares with these centres; as in part~(a) of
the proof of Lemma~\ref{lem:Kw}, the directions of the vectors $x-c_i$ of these centres are pairwise at least $\gamma_0$ apart. Let $D_i^*(I)$
be the largest number of directions, pairwise at least $\gamma_0$ apart, in which such an $x$ exists, and let $W_I$ be the number of
$c_i\in I$ that can belong to a relevant edge with a side (at most one per centre, since $\sqrt2<k-2d$). Then the number of relevant
edges is at most
\[
U(I):=E_{II}+W_I+\sum_{c_i\in I}D_i^*(I) .
\]
A branch-and-bound computation over the parameters of $I$ shows that a computed upper bound of $U(I)$ is at most $9$ on every
box of a cover of the parameter domain. Without a side,
the centres are written in polar coordinates $(r_i,\vartheta_i)$ with $\vartheta_1=0\le\vartheta_2\le\dots\le\vartheta_n$ (after a rotation and
relabelling). If a relevant edge contains a side, its rectangle $R(P,S_0)$ stands on a side of $[0,k]^2$ whose line is at
distance $h<d+\frac{\sqrt2}2$ from $p$, and the rectangles of all relevant edges with a side stand on one and the same side $\sigma$: two of them on opposite sides would force
$k<2d+\sqrt2$, which is impossible since $k\ge4$, and if $R(P,\sigma)$ and $R(P',\sigma')$ with adjacent sides $\sigma,\sigma'$ both contained $p$, then $P\ne P'$ (the rectangle of a square and a side is determined by the square) and, with the
corner at the origin, the conditions $|p_x-c_x|\le\frac12$, $0\le p_x\le c'_x<d+\frac{\sqrt2}2$, $c_x\ge\frac12$ and the symmetric ones would give
$|c-c'|^2\le\frac12$, which is impossible. We take $\sigma$ to be the line $y=-h$ and ignore the other sides (this only drops constraints); then every centre satisfies $y\ge-h+\frac12$, no normalisation of $\vartheta_1$ is used, and $h$ is a further
parameter. For a box of parameters the directions are divided into $1440$ arcs of $0.25^\circ$ and $|x-c_i|$ into $6$ intervals; an arc
is discarded only if for each interval one of the conditions $|x|>\rho_1$, $p\in R(c_i,x)$, $|x-c_j|\ge1$ (and $y\ge-h+\frac12$) fails for all
parameters in the box and all points of the arc and the interval. Then $D_i^*(I)$ is bounded by the largest number of points,
pairwise at least $\gamma_0$ apart, in the union of the remaining arcs (the largest of the circular greedy counts started at the beginning of each maximal run of
arcs, or $\lfloor2\pi/\gamma_0\rfloor$ if no arc is discarded; this gives the maximum, since an optimal set can be rotated backwards, keeping all gaps, until one of its points reaches
the start of its run, after which the $j$-th greedy point is never later than the $j$-th point of the rotated optimal set); $E_{II}$ by the number of pairs that may be at distance less
than $\Lambda$; and $W_I$ by the number of centres that may lie within $d+\frac{\sqrt2}2$ of the side with $|c_{i,x}|\le\frac12$ and $c_{i,y}\ge0$.
A box is accepted if it contains no admissible configuration (for every parameter of the box two centres are at distance less
than $1$, or a centre is closer than $\frac12$ to the side), or, in the case with a side, if no side edge is possible (then it is covered by the case without a side), or
if the bound is at most $9$; otherwise it is bisected. For every $n=1,\dots,5$, with and without a side, the computation
terminated with all boxes accepted ($30{,}365$ accepted boxes without a side and $48{,}308$ with a side).

The search itself is not relied upon: it used floating-point interval bounds with a safety margin of $10^{-9}$, and every accepted
box was re-verified in ball (midpoint-radius interval) arithmetic with rigorous error bounds, using Arb~\cite{Arb}, which has been part
of FLINT~\cite{FLINT} since version~3.0, through \texttt{python-flint}~\cite{pyflint} (version~0.9.0, with FLINT~3.6.0, at $96$-bit
working precision), boxes being bisected further where the ball enclosures were too coarse (all $78{,}673$ boxes confirmed, with $59{,}238$ further bisections and largest rigorous bound $9$).
In this re-verification the arcs of directions are closed and cover the whole circle (the floating-point end point of each arc is
replaced, where necessary, by the starting point of the next one, and the floating-point value $\widetilde{2\pi}<2\pi$ by the exact
$2\pi$), and every interval of an angle $\vartheta_i$ ending at $\widetilde{2\pi}$ is extended to the exact $2\pi$. This covers every
admissible configuration with an angle in the arc $(\widetilde{2\pi},2\pi)$, which has length less than $2.5\cdot10^{-16}$: two centres
with angles in this arc, or (without a side, where $\vartheta_1=0$) one such centre and $c_1$, would be at distance less than
$\rho_1-\frac12+10^{-15}<1$, so only $\vartheta_n$ can lie in the arc, and only in the case with a side; replacing $\vartheta_n$ by
$\widetilde{2\pi}$ then gives a configuration with ordered angles in the domain of the search, hence in an accepted box (the accepted boxes
leave no gap, see below), whose $\vartheta_n$-interval ends at $\widetilde{2\pi}$ and is therefore extended, and the extended box contains the
original configuration. The final count of
points pairwise at least $\gamma_0$ apart in a union of arcs is a finite combinatorial computation on the arc end points (the greedy
count described above); it is done in double precision, with $\gamma_0$ decreased by $10^{-12}$ and the condition for closing the circle
relaxed by a further $10^{-12}$. Each point that it considers is the start of an arc plus $j\gamma_0$ for some $0\le j\le8$, so each of its
comparisons is between $j\gamma_0$, $1\le j\le9$, and an integer multiple of the arc length $\pi/720$, and these always differ by more
than $4\cdot10^{-5}$; every number compared arises from the start of the run, the arc length, $\gamma_0-10^{-12}$ and the previous
point by a few floating-point operations on numbers of absolute value less than $14$, and lies within $10^{-14}$ of the exact quantity it
stands for. Hence every comparison has the same outcome as with the exact $\gamma_0$ and the closed arcs, and the computed count is the
exact one. The re-verification takes from the search programs the grids of arcs and of radial intervals, the floating-point value of
$\gamma_0$, the floating-point test that proposes arcs to be discarded (each discard is then certified in Arb, so a rounding error there
can only keep an arc), this final count and the bisection rule; the labels of the boxes therefore rest on this one implementation.
A separate program checks that the accepted boxes leave no gap: it rebuilds the bisection tree from the initial boxes, accepting
a node only if it coincides with a recorded box (each used exactly once) or contains no configuration with ordered angles, and
confirms in exact rational arithmetic that the volumes of these boxes add up to that of the initial boxes (this checker reuses
the bisection rule of the search, so it is not an independent implementation). The initial boxes and the radial range $[1,\Lambda]$ of the outer centres
contain the parameter domain up to the arc $(\widetilde{2\pi},2\pi)$ just mentioned, since the floating-point values used for
$\rho_1$, $\Lambda$ and $d+\frac{\sqrt2}2$ are upper bounds of the exact values (checked in high-precision arithmetic);
in the re-verification the conditions themselves are evaluated with the exact constants, the floating-point values serving only
as end points of the initial boxes.
The programs, the lists of accepted boxes and the logs are part of the accompanying material, available at
\url{https://github.com/squarepacker/k2-minus-c} and archived on Zenodo (\url{https://doi.org/10.5281/zenodo.23164301}, all versions); they were run with Python~3.12.10,
\texttt{python-flint}~0.9.0, NumPy~2.5.3 and mpmath~1.4.1. That the accepted boxes leave no gap has also been confirmed by two further checks,
written independently of the search programs and using different methods (exact rational volumes of the boxes intersected with the
region of ordered angles, and a recursive covering of the domain); they were written by separate sessions of the same AI system that
helped to write all the programs (see the section on the use of AI), and they share with the search the lists of boxes and the description
of the domain. The verification of the individual boxes (the Arb re-verification) has not yet been reproduced by an independent
implementation; the Squares Project reports a complete replay of it with these programs (from the release of version~1.0, whose programs
and data for this lemma are identical), with the same results \cite{LevR}; this repeats the same implementation.
\end{proof}

\begin{remark}
The proof bounds the number of \emph{all} pairs of squares (or of a square and one of two opposite sides) at distance less than $d$ whose
rectangles contain $p$, without using the graph $G$, and for this number the bound $9$ is sharp: the configuration described in
Remark~\ref{rem:const} below (two axis-parallel squares on both sides of a slit through $p$, each nearly touching two slightly
tilted squares on either side) has $9$ such pairs at an uncovered point; we checked this in $80$-digit arithmetic for
$d=10^{-6},10^{-5},10^{-4}$ (and beyond the range of the lemma for $10^{-3}$, $10^{-2}$), hence, by monotonicity in $d$,
for every $d\in[10^{-6},10^{-4}]$ (the coordinates and the check are part of the accompanying material). We do not know whether $9$ edges of a single graph $G_h$ or $G_v$ can meet in this way.
This remark is not used in the proofs.
\end{remark}

Apart from the definition of $R_e$, Lemma~\ref{lem:count}(b) is not used below; it bounds the multiplicity at all points,
not only at uncovered ones.

\begin{lemma}[slits]\label{lem:slit}
Let $A,B$ be distinct squares, or let one of them be the left or the right side of $[0,k]^2$, and let $E_{AB}$ be the set
of $y\in[0,k]$ such that on the line $y$ both $A$ and $B$ are met, $A$ lies to the left of $B$, and no square is met
strictly between them. For $y\in E_{AB}$ let $g(y)$ be the length of the gap between them, and let $\theta:=\theta(A,B)$.
\begin{enumerate}[label=(\alph*)]
\item All $y\in E_{AB}$ with $g(y)<1$ lie in one component of $E_{AB}$.
\item On every component $J$ of $E_{AB}$ the function $g$ is convex, and $\int_Ig\ge\theta(2-\theta)|I|^2/4$ for every interval
$I\subset J$. In particular, for $\eta>0$ the set $\{y\in J:g(y)<\eta\}$ is an interval.
\end{enumerate}
The same holds with the roles of the axes exchanged (vertical lines, $A$ below $B$, the bottom and the top side).
\end{lemma}

\begin{proof}
Let $U$ be the intersection of the vertical extents of $A$ and $B$ (for a side of the container, its vertical extent is
$[0,k]$); it is an interval containing $E_{AB}$. Let $r_A(y)$ be the largest abscissa of $A$ on the line $y$ ($r_A\equiv0$ for the left
side) and $l_B(y)$ the smallest abscissa of $B$ ($l_B\equiv k$ for the right side); by Lemma~\ref{lem:bd} (with $(s,t)=(x,y)$),
$G:=l_B-r_A$ is convex on $U$, and $g=G$ on $E_{AB}$.
If a square $Z$ lies strictly between $A$ and $B$ on one line
$y\in U$, then it does so on every line meeting $A$, $B$ and $Z$: two disjoint compact convex sets are strictly
separated by a line, and if that line is not horizontal their order on every horizontal line meeting both is the same (if
it is horizontal they never share a line); when one of $A,B$ is a side of the container this is trivial, since every square
lies to the right of the left side and to the left of the right side.

(a) Let $y_1<y_2$ be points of $E_{AB}$ with $g(y_1),g(y_2)<1$; by convexity $G<1$ on $[y_1,y_2]$. Suppose that some
$y\in(y_1,y_2)$ is not in $E_{AB}$. Then a square $Z$ is met strictly between $A$ and $B$ on the line $y$; it does not meet the
lines $y_1$ and $y_2$, since it would lie between $A$ and $B$ there, so its vertical extent $V(Z)$ lies in $(y_1,y_2)\subset U$.
The midpoint $y_Z$ of $V(Z)$ lies in the middle band $\mathcal M(Z)$, so the chord of $Z$ at height $y_Z$ has length
$\sec a(Z)\ge1$ (Lemma~\ref{lem:geom}(b)), and it lies strictly between $A$ and $B$. Hence $G(y_Z)>1$, a contradiction. So
$[y_1,y_2]\subset E_{AB}$.

(b) On a component $J$ of $E_{AB}$, $g=G$ is convex. By Lemma~\ref{lem:bd}, it is piecewise linear with at most two break points,
at each of which its slope increases by at least $2$, and off the break points $|g'|=|\tan\psi_B-\tan\psi_A|\ge|\psi_B-\psi_A|\ge\theta$,
where $\psi_A,\psi_B$ are the angles between the vertical and the sides of $A$ and $B$ bounding the gap (for a side of the container
the angle is $0$, and $\theta(A,S_0)$ is the inclination of the square). As $\theta\le\frac\pi4<1$, Lemma~\ref{lem:cg} applies
on every interval $I\subset J$. Sublevel sets of a convex function on an interval are intervals.
\end{proof}

\begin{lemma}[cost of a large inclination on a line]\label{lem:Bh}
Let $0<\bar\beta\le1.5\cdot10^{-6}$, and let $Y\subset[0,k]$ and $\beta:Y\to[0,\bar\beta]$ be such that for every $y\in Y$:
$\omega(y)<\frac12$ and the line $y$ meets a square of inclination $\ge\beta(y)$. Then
\[
\int_Y2\beta(y)\,dy\ \le\ A\,W,\qquad A:=\frac{4\sqrt{K_w^*/Q_*}}{2-3\cdot10^{-6}}\ <\ 10.6067 .
\]
\end{lemma}

\begin{proof}
For $y\in Y$ let $X_1,\dots,X_m$ be the squares met by the line, in order, and $X_0,X_{m+1}$ the left and right sides of
the container. Each gap is part of the line waste, so it is shorter than $\frac12$, and the triangle
inequality from the left side to a square of inclination $\ge\beta(y)$ and on to the right side gives
$\sum_{j=0}^m\theta(X_j,X_{j+1})\ge2\beta(y)$. Put $\theta_{\max}:=2\bar\beta$ and $\theta':=\min(\theta,\theta_{\max})$; then also
$\sum_j\theta'(X_j,X_{j+1})\ge2\beta(y)$, since either some term equals $\theta_{\max}\ge2\beta(y)$, or no term is capped.

For a pair $e=\{P,Q\}$ of members of $\mathcal V_h$, with $P$ the left one,
let $Y_e$ be the set of $y\in Y$ on which $P$ and $Q$ are consecutive, and let $\mathcal E$ be the finite set of pairs with
$|Y_e|>0$; they are edges of $G_h$. Write $\theta_e:=\theta(P,Q)$, $\theta_e':=\min(\theta_e,\theta_{\max})$, $E_e:=E_{PQ}\supset Y_e$, and
$g_e$ for the gap function of Lemma~\ref{lem:slit}. Integrating over $Y$,
\[
\int_Y2\beta\ \le\ \sum_{e\in\mathcal E}\theta_e'\,|Y_e| .
\]
Put $d':=\theta_{\max}$ and split $Y_e=Y_e^c\cup Y_e^f$, where $Y_e^c:=\{y\in Y_e:g_e(y)<d'\}$. Let
$s_e^f:=\int_{E_e\cap\{g_e\ge d'\}}g_e$. Since $\theta_e'\le\theta_{\max}=d'$,
\[
\theta_e'|Y_e^f|\ \le\ \frac{\theta_{\max}}{d'}\int_{Y_e^f}g_e\ \le\ s_e^f .
\]
Now let $|Y_e^c|>0$; then the members of $e$ are at distance $<d'\le3\cdot10^{-6}$. As $g_e<\frac12$ on $Y_e$, Lemma~\ref{lem:slit}
shows that $Y_e$ lies in one component $J_e$ of $E_e$, and that $Y_e^c$ lies in the interval $I_e:=\{y\in J_e:g_e(y)<d'\}$, with
\[
s_e^c:=\int_{I_e}g_e\ \ge\ \frac{\theta_e(2-\theta_e)}4|I_e|^2\ \ge\ \frac{\theta_e'(2-\theta_{\max})}4|I_e|^2 ,
\]
because $x(2-x)$ is increasing on $[0,1]$. Let $f_e$ be the area of the waste in $R_e$ (notation of
Lemma~\ref{lem:count}(b), with $d=d'$). We claim that $f_e\ge Q_*\theta_e'(2-\theta_{\max})/4$. If $\theta_e\le\theta_1:=10^{-5}$, this
follows from Proposition~\ref{prop:FLpar}. Otherwise Lemma~\ref{lem:count}(a) gives
$f_e\ge\theta_e(2-\theta_e)/(4q_0)\ge\theta_1(2-\theta_1)/40$, which exceeds $Q_*\theta_{\max}(2-\theta_{\max})/4\ge Q_*\theta_e'(2-\theta_{\max})/4$ as
$\theta_{\max}\le3\cdot10^{-6}$. (For a square $X$ and a side, $R_e=R(X,S_0)$ may stand on another side of the container; this is
harmless, as $\theta(X,S_0)=a(X)$.) Hence
\[
\theta_e'|Y_e^c|\ \le\ \theta_e'|I_e|\ =\ \big(\theta_e'|I_e|^2\cdot\theta_e'\big)^{1/2}\ \le\ \frac4{2-\theta_{\max}}\Big(\frac{s_e^c\,f_e}{Q_*}\Big)^{1/2}.
\]
The sets $\{(x,y):y\in E_e,\ x\text{ in the gap of }e\text{ on the line }y\}$ are disjoint parts of the waste for different
$e$, and $I_e$ and $E_e\cap\{g_e\ge d'\}$ are disjoint, so $S^c+S^f\le W$ for $S^c:=\sum_es_e^c$ and $S^f:=\sum_es_e^f$. By
Lemma~\ref{lem:Kw9} (with $d=d'$), every point of the waste lies in $R_e$ for at most $K_w^*$ of the pairs with $|Y_e^c|>0$, so
$\sum f_e\le K_w^*W$, the sum over these pairs. By the Cauchy--Schwarz inequality,
\[
\sum_{e\in\mathcal E}\theta_e'|Y_e|\ \le\ \frac{4\sqrt{K_w^*/Q_*}}{2-\theta_{\max}}\sqrt{S^cW}+S^f\ \le\ A\sqrt{(W-S^f)W}+S^f\ \le\ AW,
\]
the last step because $x\mapsto A\sqrt{(W-x)W}+x$ is decreasing on $[0,W]$ when $A\ge2$. Finally $4\sqrt{9/0.32}/(2-3\cdot10^{-6})<10.6067$.
\end{proof}

\begin{lemma}[cost of a large inclination in a column]\label{lem:Bv}
Let $0<\delta\le10^{-5}$ and $0<\bar\beta\le3\cdot10^{-6}$. Let $X,X'\subset[0,k]$. For $x\in X$ let $\bar y_x\in(0,\frac k2-1]$ and
$\beta(x)\in[0,\bar\beta]$ be such that the segment $\sigma_x=\{x\}\times[0,\bar y_x]$ has vertical waste $<\delta$ and meets a square of
inclination $\ge\beta(x)$; for $x\in X'$ let $\bar y'_x$, $\beta'(x)$ satisfy the same with the segment
$\sigma_x'=\{x\}\times[k-\bar y'_x,k]$. Then
\[
\int_X\beta(x)\,dx+\int_{X'}\beta'(x)\,dx\ \le\ A\,W .
\]
\end{lemma}

\begin{proof}
For $x\in X$ list the squares met by $\sigma_x$ from the bottom up to the first one of inclination $\ge\beta(x)$, preceded by the
bottom side of the container; for $x\in X'$ list those met by $\sigma'_x$ from the top down, preceded by the top side. All gaps
are $<\delta$, so consecutive members are at distance $<\delta$, and the triangle inequality gives
$\sum_j\theta(Y_j,Y_{j+1})\ge\beta(x)$ (resp.\ $\beta'(x)$), hence the same for $\theta':=\min(\theta,\bar\beta)$. Two members that are
consecutive on $\sigma_x$ are consecutive on the whole vertical line through $x$ (a square between them would meet $\sigma_x$
between them), so the pairs used for a set of $x$ of positive measure are edges of $G_v$ whose members are at distance $<\delta$.
Both members of a pair used for some $x\in X$ meet $\R\times[0,\frac k2-1]$, and both members of a pair used for some $x\in X'$ meet
$\R\times[\frac k2+1,k]$; as no square meets both strips, no pair is used in both ways. For a pair $e$ let $X_e$ be the set of
$x$ in $X$, or in $X'$, at which it is used. Since the gaps are $<\delta<1$, the vertical version of Lemma~\ref{lem:slit} puts $X_e$
in one component $J_e$ of $E_e$, and $s_e:=\int_{J_e}g_e\ge\theta_e'(2-\bar\beta)|J_e|^2/4$. As in the proof of Lemma~\ref{lem:Bh}, with
$d=\delta$ and $\theta_{\max}$ replaced by $\bar\beta$, we have $f_e\ge Q_*\theta_e'(2-\bar\beta)/4$ (by Proposition~\ref{prop:FLpar} if
$\theta_e\le\theta_1=10^{-5}$, and by Lemma~\ref{lem:count}(a) otherwise, since $\theta_1(2-\theta_1)/40\ge Q_*\bar\beta(2-\bar\beta)/4$). Hence
$\theta_e'|X_e|\le\theta_e'|J_e|\le\frac4{2-\bar\beta}(s_ef_e/Q_*)^{1/2}$, $\sum_es_e\le W$ and $\sum_ef_e\le K_w^*W$ (Lemma~\ref{lem:Kw9} with $d=\delta$), so the
left-hand side is at most $\sum_e\theta_e'|X_e|\le\frac{4\sqrt{K_w^*/Q_*}}{2-\bar\beta}W\le AW$ (as $\bar\beta\le3\cdot10^{-6}$).
\end{proof}

For $y\in(0,k)$ let $d(y):=\min(y,k-y)$ be the distance to the nearer of the bottom and the top side.

\begin{lemma}[excess of long chords on a line]\label{lem:Eline}
Let $0<\delta\le10^{-5}$ and $0<\alpha\le1.5\cdot10^{-6}$, let $y\in(0,k)$ with $d(y)\le\frac k2-1$, and suppose that every square
meeting the line $y$ has inclination $<\alpha$. Then, with $E(y)$ as in Lemma~\ref{lem:E},
\[
E(y)\ \le\ 0.50001\,\alpha\Big(A+\frac\alpha\delta\Big)W .
\]
\end{lemma}

\begin{proof}
We use $\sec x-1\le0.50001\,x^2$ for $0\le x\le10^{-3}$ (the function $(\sec x-1)/x^2$ has a Taylor series with positive
coefficients, so it is increasing on $(0,\frac\pi2)$, and its value at $10^{-3}$ is less than $0.50001$). Let $S$ be a square with
$c_S(y)\ge1$, and $a:=a(S)<\alpha$. By Lemma~\ref{lem:geom}(b), $c_S(y)\le\sec a$, so
\[
c_S(y)-1\ \le\ 0.50001\,a^2\ \le\ 0.50001\,\alpha\,a\,c_S(y),
\]
as $a<\alpha$ and $c_S(y)\ge1$. The chords of these squares on the line $y$ are disjoint; let $U\subset[0,k]$ be the union of their
projections on the $x$-axis and, for $x\in U$, let $S_x$ be the square whose chord contains $(x,y)$. Summing,
\[
E(y)\ \le\ 0.50001\,\alpha\sum_{S:\,c_S(y)\ge1}a(S)\,c_S(y)\ =\ 0.50001\,\alpha\int_Ua(S_x)\,dx .
\]
Let $y\le\frac k2$; then $0<y=d(y)\le\frac k2-1$. For $x\in U$ the segment $\sigma_x:=\{x\}\times[0,y]$ meets $S_x$, of inclination
$a(S_x)<\alpha$. Let $U'$ be the set of $x\in U$ for which $\sigma_x$ has vertical waste $\ge\delta$; as the vertical waste of
$\sigma_x$ is at most that of $\{x\}\times[0,k]$, whose integral over $x$ is $W$, we have $|U'|\le W/\delta$, and so
$\int_{U'}a(S_x)\,dx\le\alpha W/\delta$. Lemma~\ref{lem:Bv}, applied to $X:=U\setminus U'$, $X':=\emptyset$, $\bar y_x:=y$, $\beta(x):=a(S_x)$
and $\bar\beta:=\alpha$, gives $\int_{U\setminus U'}a(S_x)\,dx\le AW$. If $y>\frac k2$, the same argument with the segments
$\{x\}\times[y,k]$ hanging from the top side (that is, $X:=\emptyset$, $X':=U\setminus U'$ and $\bar y'_x:=k-y=d(y)$) gives the same bound.
Hence $\int_Ua(S_x)\,dx\le(A+\alpha/\delta)W$.
\end{proof}

\section{Quantization in columns}

Fix $C:=2$ and $\delta:=10^{-5}$, and put $\alpha(s):=1/(Cs)$ for $s>0$.

\begin{definition}
Let $0<s\le\frac k2-3$. A point $x\in[0,k]$ is \emph{$s$-rigid from below} if the segment
$\sigma^s_x:=\{x\}\times[0,s+2]$ has vertical waste $<\delta$ and meets no square of inclination $\ge\alpha(s)$, and
\emph{$s$-rigid from above} if the segment $\tilde\sigma^s_x:=\{x\}\times[k-s-2,k]$ has these two properties.
We write $\mathrm{NR}_s$ and $\widetilde{\mathrm{NR}}_s$ for the sets of $x$ that are not $s$-rigid from below, respectively
from above.
\end{definition}

If $s<s'$, then $\sigma^s_x\subset\sigma^{s'}_x$ and $\alpha(s)>\alpha(s')$, so $\mathrm{NR}_s\subset\mathrm{NR}_{s'}$. Moreover, for fixed $x$
the set of $s\in(0,\frac k2-3]$ with $x\in\mathrm{NR}_s$ is closed: it is the union of the set where the vertical waste of
$\sigma^s_x$, a continuous function of $s$, is $\ge\delta$, and of the sets $\{s:\ s+2\ge b_X,\ Cs\,a(X)\ge1\}$ over the finitely many
squares $X$ with $a(X)>0$ that meet the vertical line through $x$, where $b_X$ is the lower end point of the chord of $X$ on
that line. The same holds for $\widetilde{\mathrm{NR}}_s$.

For a square $Z$ and $x\in\mathcal M'(Z)$ let $\beta_Z(x)$ be the height of the lower endpoint of the vertical chord of
$Z$ at $x$. A square $X\ne Z$ is \emph{below $Z$ at $x$} if $X$ meets $\{x\}\times[0,\beta_Z(x))$; its vertical chord at
$x$ is then contained in $[0,\beta_Z(x))$, because it is an interval disjoint from the chord of $Z$.

\begin{lemma}[quantization]\label{lem:Q}
Let $s\ge2.36\cdot10^5$ with $s\le\frac k2-3$, let $Z$ be a square with $a(Z)<\alpha(s)$ whose lowest point has height
$v_Z\le s$, and let $x\in\mathcal M'(Z)$ satisfy
\begin{enumerate}[label=(\roman*)]
\item $x\in\mathcal M'(X)$ for every square $X$ with $a(X)<\alpha(s)$ that is below $Z$ at $x$, and
\item $x$ is $s$-rigid from below.
\end{enumerate}
Then every point of $\mathcal R(Z)$ is at distance less than $w_0:=1.265\cdot10^{-5}$ from an integer.
\end{lemma}

\begin{proof}
Write $\beta=\beta_Z(x)$ and $\alpha=\alpha(s)\le2.119\cdot10^{-6}$; by Lemma~\ref{lem:geom}(c), $\beta\le v_Z+\sin a(Z)<s+\alpha$,
so $\{x\}\times[0,\beta)\subset\sigma^s_x$. By (ii) every square below $Z$ at $x$ has inclination $<\alpha$, so by (i) and
Lemma~\ref{lem:geom}(c) its vertical chord at $x$ has length in $[1,\sec\alpha]$. If $n$ is the number of these squares,
their chords and the uncovered part fill $[0,\beta)$, whence $n\le\beta<n\sec\alpha+\delta$ and $n<s+\alpha$. We use
$\sec x-1\le0.50001\,x^2$ for $0\le x\le10^{-3}$ (see the proof of Lemma~\ref{lem:Eline}). Now
$n(\sec\alpha-1)\le(s+\alpha)\cdot0.50001\,\alpha^2=0.50001(1+\alpha/s)/(C^2s)<5.3\cdot10^{-7}$, so $\beta\in[n,\ n+1.053\cdot10^{-5})$. By
Lemma~\ref{lem:geom}(c), and as $\sin a(Z)<2.119\cdot10^{-6}$, $v_Z\in[\beta-\sin a(Z),\beta]\subset(n-2.119\cdot10^{-6},\ n+1.053\cdot10^{-5})$.
The lower ramp of $Z$ lies in $[v_Z,v_Z+\sin a(Z)]\subset(n-2.119\cdot10^{-6},\ n+1.2649\cdot10^{-5})$. The upper ramp is
$[v_Z+\cos a(Z),\ v_Z+\cos a(Z)+\sin a(Z)]$; here $v_Z+\cos a(Z)\ge\beta-\sin a(Z)+\cos a(Z)\ge n+\cos\alpha-\sin\alpha
>n+1-2.12\cdot10^{-6}$ and $v_Z+\cos a(Z)+\sin a(Z)\le\beta+1+\sin a(Z)<n+1+1.2649\cdot10^{-5}$.
\end{proof}

For a square $X$ with $a(X)>0$ and horizontal extent $[x_L,x_R]$, the vertical ramp set of $X$ (Lemma~\ref{lem:geom}(c))
consists of two intervals of length $\sin a(X)$: the \emph{left end} $[x_L,x_L+\sin a(X)]$ and the \emph{right end}
$[x_R-\sin a(X),x_R]$. They are the projections on the $x$-axis of the \emph{left side} and the \emph{right side} of $X$, the sides
through the leftmost and the rightmost vertex that make the angle $a(X)$ with the vertical. If $a(X)=0$ we say that the left
end of $X$ is the single point $x_L$ (the abscissa of its left side).

\begin{lemma}[cracks]\label{lem:crack}
Let $0<\alpha\le10^{-3}$ and $0<\delta\le10^{-3}$ \textup{(}in this lemma $\alpha$ and $\delta$ are arbitrary numbers in these ranges\textup{)}, and let $\sigma=\{x\}\times[0,b]$ be a vertical segment with vertical waste
$<\delta$ such that every square meeting $\sigma$ has inclination $<\alpha$. Let $X$ be a square meeting $\sigma$ such that $x$ lies in
the right end $[\xi_X,\xi_X+\sin a(X)]$ of $X$ \textup{(}so $0<a(X)<\alpha$\textup{)}, let $p=(x,h)$ be the point where the vertical
line through $x$ meets the right side of $X$, and assume $h\le b-\delta$. Then either
$x-\xi_X<\delta\tan\alpha\,(1+2.01\tan^2\alpha)$, or there is a square $X'\ne X$ that meets $\sigma$ at a point $q\ne p$ of the
vertical segment of length $\delta$ that starts at $p$ and leaves $X$ there \textup{(}so $q\notin X$ and $|q-p|<\delta$\textup{)}, and satisfies
\[
\lambda'\ \le\ x\ \le\ \lambda''+2.01\,\delta\tan\alpha ,
\]
where $[\lambda',\lambda'']$ is the left end of $X'$.
\end{lemma}

\begin{proof}
Write $a:=a(X)$, let $\varsigma$ be the right side of $X$, $P_1$ its end point at the rightmost vertex (abscissa $\xi_X+\sin a$), and $P$
its other end point (abscissa $\xi_X$); let $P_y$ denote the height of $P$. Then $\varsigma$ runs from $P_1$ in the direction $(-\sin a,u\cos a)$ for some $u\in\{1,-1\}$, and
for heights $y$ between those of $P_1$ and $P$ the right end point of the chord of $X$ at height $y$ is $r(y):=x-u(y-h)\tan a$.
Suppose that $x-\xi_X\ge\delta\tan\alpha\,(1+2.01\tan^2\alpha)$. Then $|P_y-h|=(x-\xi_X)\cot a\ge\delta(1+2.01\tan^2\alpha)$, so the
heights $h+ut$ with $0\le t\le\delta(1+2.01\tan^2\alpha)$ lie between those of $p$ and $P$. Moreover the segment
$\{x\}\times\{h+ut:0<t<\delta\}$ lies in $\sigma$: for $u=1$ because $h\le b-\delta$, for $u=-1$ because $P_y\ge0$. As its uncovered
part has length $<\delta$, it contains a point $q=(x,y_q)$, $y_q=h+ut$, $0<t<\delta$, of some square $X'$. Since
$r(y_q)=x-t\tan a<x$, the point $q$ is not in $X$, so $X'\ne X$, and $a':=a(X')<\alpha$ because $X'$ meets $\sigma$.

The disjoint convex sets $X$ and $X'$ both meet the horizontal line through $q$, so a line separating them is not horizontal, and
on every horizontal line meeting both, $X$ lies to the left of $X'$. Let $l'(y)$ be the left end point of the chord of $X'$ at
height $y$. Then (A) $l'(y)>r(y)$ for every height $y$ between those of $p$ and $P$ at which $X'$ is met, and (B)
$x-t\tan a=r(y_q)<l'(y_q)\le x$. Let $\hat q:=(l'(y_q),y_q)$. If $a'=0$, then $\hat q$ lies on the vertical left side of $X'$, and (B)
gives the claim with $\lambda'=\lambda''=l'(y_q)$. Let $a'>0$, and let $Q'$ be the leftmost vertex of $X'$, so that
$\lambda'=Q'_x$. The left boundary $\{(l'(y),y)\}$ of $X'$ consists of the left side, which projects onto $[\lambda',\lambda'']$,
and of the other side $F'$ of $X'$ at $Q'$, which runs from $Q'$ in the direction $(\cos a',v\sin a')$ for some $v\in\{1,-1\}$.

\emph{Case 1: $\hat q$ lies on the left side.} Then $l'(y_q)\in[\lambda',\lambda'']$, and by (B) $\lambda'\le x<\lambda''+\delta\tan\alpha$.

\emph{Case 2: $\hat q=Q'+\mu(\cos a',v\sin a')\in F'$ with $0<\mu\le1$, and $v=-u$.} For $0<w\le\mu$ the point
$\hat q-w(\cos a',v\sin a')$ of $X'$ lies at height $y_q+uw\sin a'$. If this height lies between those of $p$ and $P$, then (A)
gives $l'(y_q)-w\cos a'>r(y_q+uw\sin a')=x-(t+w\sin a')\tan a$, that is, $w(\cos a'-\sin a'\tan a)<l'(y_q)-x+t\tan a\le t\tan a$
by (B). Put $w_*:=t\tan a/(\cos a'-\sin a'\tan a)$. If $\mu>w_*$, applying this to all $w\in(w_*,\mu]$ shows that $|P_y-y_q|\le
w_*\sin a'$, hence $|P_y-h|\le t+w_*\sin a'<\delta(1+2.01\tan^2\alpha)$, contrary to the first paragraph. So $\mu\le w_*$, and by (B)
$x-\lambda'=(x-l'(y_q))+\mu\cos a'<t\tan a+t\tan a/(1-\tan a\tan a')<2.01\,\delta\tan\alpha$, while $x\ge l'(y_q)\ge\lambda'$.

\emph{Case 3: $\hat q=Q'+\mu(\cos a',v\sin a')\in F'$ with $0<\mu\le1$, and $v=u$.} Then the left side of $X'$ runs from $Q'$ in the
direction $(\sin a',-u\cos a')$, and $Q'_y=y_q-u\mu\sin a'$, while $h=y_q-ut$. If $\mu\sin a'>t$, the height $h$ lies strictly between
$Q'_y$ and $y_q$, so $l'(h)=l'(y_q)-t\cot a'$; by (A) and (B), $l'(h)>r(h)=x\ge l'(y_q)$, which is absurd. So $\mu\sin a'\le t$; then
the height $h$ lies on the left side of $X'$, at distance $t-\mu\sin a'<\delta$ from $Q'_y$, so that $l'(h)\le\lambda'+\delta\tan\alpha$,
and (A) gives $x=r(h)<l'(h)\le\lambda'+\delta\tan\alpha$; also $x\ge l'(y_q)\ge\lambda'$.

In all cases $\lambda'\le x\le\lambda''+2.01\,\delta\tan\alpha$, and $q\in X'\cap\sigma$ is at distance $t<\delta$ from $p$.
\end{proof}

\begin{lemma}[room in a column]\label{lem:GZ}
Under the assumptions of Lemma~\ref{lem:Q} on $s$ and $Z$, let $B_Z=B_Z^{(s)}$ be the set of points $x\in\mathcal M'(Z)$ that are
$s$-rigid from below but violate~(i). Then $|B_Z|<0.500023$, so that $G_Z=G_Z^{(s)}:=\mathcal M'(Z)\setminus B_Z$ has measure greater than
$g_0:=0.49997$.
\end{lemma}

\begin{proof}
Write $\alpha=\alpha(s)\le2.119\cdot10^{-6}$, $\varrho:=\cos\alpha+\sin\alpha$ and $\mathcal M'(Z)=[m_0,m_1]$, so that
$m_1-m_0=\cos a(Z)-\sin a(Z)\le1$; note that the left end of $Z$ ends at $m_0$ (if $a(Z)=0$ it is the point $m_0$; if $a(Z)>0$,
then in the coordinates of the proof of Lemma~\ref{lem:geom} the horizontal extent of $Z$ is $[x_0-\sin a(Z),x_0+\cos a(Z)]$, its left end
is $[x_0-\sin a(Z),x_0]$ and $\mathcal M'(Z)=[x_0,x_0+\cos a(Z)-\sin a(Z)]$, and likewise in the reflected position). Let $\mathcal X_{\mathrm L}$ be the set of squares $X$ with
$a(X)<\alpha$ that are below $Z$ at some point of $\mathcal M'(Z)$ and whose left end is contained in $(m_0,m_1)$, and let
$\mathcal X_{\mathrm R}$ be the set of squares $X$ with $0<a(X)<\alpha$ that are below $Z$ at some point of $\mathcal M'(Z)$ and whose right end
is contained in $(m_0,m_1)$.

\emph{Counting.} If $X\in\mathcal X_{\mathrm L}$ and $y\in\mathcal M(X)$ \textup{(}every $y$ in the vertical extent of $X$ if $a(X)=0$\textup{)}, the
chord of $X$ at height $y$ has its left end point in the left end of $X$, hence in $(m_0,m_1)$, and length $\ge1\ge m_1-m_0$, so it
contains $m_1$. Chords of distinct squares at the same height are disjoint, so the middle bands of the squares in $\mathcal X_{\mathrm L}$
have pairwise disjoint interiors. A square $X$ below $Z$ at some $x\in\mathcal M'(Z)$ meets $\{x\}\times[0,\beta_Z(x))$, and
$\beta_Z(x)\le v_Z+\sin a(Z)<s+\alpha$; so its middle band lies in $[0,s+\alpha+\varrho]$ and has length $\cos a(X)-\sin a(X)>
\cos\alpha-\sin\alpha$. Hence $|\mathcal X_{\mathrm L}|\le n_*:=(s+\alpha+\varrho)/(\cos\alpha-\sin\alpha)$, and in the same way (with $m_0$ in
place of $m_1$) $|\mathcal X_{\mathrm R}|\le n_*$.

\emph{Covering.} Let $x\in B_Z$. As (i) fails, some square $X$ with $a(X)<\alpha$ is below $Z$ at $x$ and $x\notin\mathcal M'(X)$; then
$a(X)>0$ and $x$ lies in the left or in the right end of $X$. If that end is not contained in $(m_0,m_1)$, it contains $m_0$ or
$m_1$, and $x\in[m_0,m_0+\alpha)\cup(m_1-\alpha,m_1]$. If it is the left end and it is contained in $(m_0,m_1)$, then $X\in\mathcal X_{\mathrm L}$.
If it is the right end and it is contained in $(m_0,m_1)$, then $X\in\mathcal X_{\mathrm R}$, and we apply Lemma~\ref{lem:crack} to
$\sigma=\sigma^s_x$ (here $b=s+2$ and $h<\beta_Z(x)<s+\alpha\le b-\delta$; all squares meeting $\sigma^s_x$ have inclination $<\alpha$, as $x$ is
rigid). Either $x\in[\xi_X,\xi_X+1.01\,\delta\tan\alpha]$, or there is a square $X'$ as in the lemma; let $q$ be the point of $X'\cap\sigma$
found there. If $X'=Z$, then $x\le m_0+2.01\,\delta\tan\alpha$. Otherwise $q\notin Z$, and since $q$ lies within $\delta$ of $p$, which is below
the chord $[\beta_Z(x),\beta_Z(x)+\sec a(Z)]$ of $Z$, the point $q$ lies below that chord; so $X'$ is below $Z$ at $x$, and
$a(X')<\alpha$. If the left end of $X'$ is contained in $(m_0,m_1)$, then $X'\in\mathcal X_{\mathrm L}$ and $x\in[\lambda',\lambda''+2.01\,\delta\tan\alpha]$;
otherwise $\lambda'\le m_0$ or $\lambda''\ge m_1$; in the first case $\lambda''<m_0+\alpha$ and $x<m_0+\alpha+2.01\,\delta\tan\alpha$, and in
the second case $x\ge\lambda'>\lambda''-\alpha\ge m_1-\alpha$ (if $a(X')=0$, then $x\ge\lambda'=\lambda''\ge m_1$). Altogether,
\[
B_Z\ \subset\ [m_0,m_0+1.01\alpha]\cup[m_1-\alpha,m_1]\cup\bigcup_{X\in\mathcal X_{\mathrm L}}[\lambda'_X,\lambda''_X+2.01\,\delta\tan\alpha]\cup
\bigcup_{X\in\mathcal X_{\mathrm R}}[\xi_X,\xi_X+1.01\,\delta\tan\alpha].
\]
As $\lambda''_X-\lambda'_X=\sin a(X)<\alpha$ and $\tan\alpha<1.000001\,\alpha$,
\begin{align*}
|B_Z|\ &\le\ 2.01\,\alpha+n_*\alpha\,(1+3.02\,\delta\cdot1.000001)\\
&\le\ 2.01\,\alpha+\frac{(1+(\alpha+\varrho)/s)(1+3.04\cdot10^{-5})}{C(\cos\alpha-\sin\alpha)}\ <\ 0.500023,
\end{align*}
since $\alpha=1/(Cs)\le2.119\cdot10^{-6}$, $s\ge2.36\cdot10^5$, $\delta=10^{-5}$ and $C=2$ (the middle expression decreases as $s$ grows); and
$|\mathcal M'(Z)|=\cos a(Z)-\sin a(Z)>1-2.12\cdot10^{-6}$.
\end{proof}

\noindent\emph{Columns from the top.} The reflection $(x,y)\mapsto(x,k-y)$ maps $[0,k]^2$ onto itself and closed packings to closed
packings; since $k$ is an integer, it maps points at distance less than $w_0$ from an integer to such points. Applying
Lemmas~\ref{lem:Q}--\ref{lem:GZ} to the reflected packing, we obtain the same statements with ``lowest point'' replaced by
``highest point'', $v_Z\le s$ by $v_Z'\ge k-s$ for the height $v_Z'$ of the highest point of $Z$, ``below'' by ``above'' (with the upper end
point of the vertical chord of $Z$), and ``rigid from below'' by ``rigid from above''; we denote the corresponding set by
$\widetilde G_Z^{(s)}$.

\section{Proof of Theorem~\ref{thm:main}}

Put $\varepsilon:=2\cdot10^{-3}$, $\omega_0:=2\cdot10^{-4}$, $y_{\min}:=236000$, $w_0:=1.265\cdot10^{-5}$ and $h_0:=1-2w_0$, and
\[
F(k):=1.99954\,(\log k-13.06675)/30.418 .
\]
For $2\le k<10^{12}$ we have $\log k<27.632$, so $F(k)<0.958<1\le k^2-N$ by Corollary~\ref{cor:W1}. Let now
$k\ge10^{12}$, and put $y_1:=\frac k2-3$; recall $C=2$, $\delta=10^{-5}$ and $\alpha(s)=1/(Cs)$ from Section~5, and put
$C':=2\sqrt2$ and $\alpha'(y):=1/(C'y)$ (the threshold for lines). Recall that $d(y)=\min(y,k-y)$ is the distance to the
nearer of the bottom and the top side, and for a measurable set $Y\subset\{y:d(y)\ge y_{\min}\}$ write $\ell(Y):=\int_Ydy/d(y)$. Let
\[
H:=\{y\in(0,k):\ y_{\min}\le d(y)\le(1-\varepsilon)y_1,\ \frc{y}\in[w_0,1-w_0]\},
\]
$H_b:=H\cap(0,\frac k2)$ and $H_t:=H\cap(\frac k2,k)$.
Since $k$ is an integer, $y\mapsto k-y$ maps $H_b$ onto $H_t$ and preserves $d$, so $\ell(H)=2\ell(H_b)$. If
$j_0:=y_{\min}\le m\le j_1-1$ with $j_1:=\lfloor(1-\varepsilon)y_1\rfloor$ (note $j_1>j_0$), then $[m+w_0,m+1-w_0]\subset H_b$ and
$\int_{m+w_0}^{m+1-w_0}dy/y\ge h_0/(m+1)\ge h_0\int_{m+1}^{m+2}dy/y$, so
\[
\ell(H)\ \ge\ 2h_0\log\frac{j_1+1}{j_0+1}\ \ge\ 2h_0\log\frac{(1-\varepsilon)(k/2-3)}{y_{\min}+1}\ \ge\ 2h_0(\log k-13.06675),
\]
the last step because $(1-\varepsilon)(k/2-3)=\frac{1-\varepsilon}2k(1-6/k)$, $\log\frac2{1-\varepsilon}+\log(y_{\min}+1)<13.066741$ and, for
$k\ge10^{12}$, $\log(1-6/k)>-10^{-11}$.

For $y\in H$ we have $\alpha'(d(y))\le\alpha'(y_{\min})<1.4982\cdot10^{-6}$ and $d(y)\le(1-\varepsilon)y_1<\frac k2-1$.
Every $y\in H$ is of exactly one of the following kinds: (i) $\omega(y)\ge\omega_0$; (ii) $\omega(y)<\omega_0$ and the line $y$ meets a
square of inclination $\ge\alpha'(d(y))$; (iii) $\omega(y)<\omega_0$ and every square meeting the line $y$ has inclination
$<\alpha'(d(y))$. Let $Y_i$, $Y_{ii}$, $Y_{iii}$ be the corresponding sets.

\smallskip\noindent\emph{Kinds (i) and (ii).} Since $\omega\ge\omega_0$ on $Y_i$ and $d\ge y_{\min}$ on $H$,
$W\ge\omega_0|Y_i|\ge\omega_0y_{\min}\,\ell(Y_i)$. Lemma~\ref{lem:Bh}, applied to $Y=Y_{ii}$ with $\beta(y)=\alpha'(d(y))\le\bar\beta:=1.4982\cdot10^{-6}$,
gives $AW\ge\int_{Y_{ii}}2\alpha'(d(y))\,dy=\frac2{C'}\ell(Y_{ii})=\frac1{\sqrt2}\ell(Y_{ii})$.

\smallskip\noindent\emph{Kind (iii): the excess of long chords.} For $y\in Y_{iii}$, Lemma~\ref{lem:Eline} with $\alpha=\alpha'(d(y))$ gives
\[
E(y)\ \le\ \gamma(y)\,W,\qquad \gamma(y):=0.50001\,\alpha'(d(y))\Big(A+\frac{\alpha'(d(y))}\delta\Big),
\]
and, bounding the parts of $H$ below and above $\frac k2$ separately (on both, $d(y)\ge y_{\min}$),
\begin{align*}
\Gamma:=\int_H\gamma(y)\,\frac{dy}{d(y)}\ &\le\ 2\cdot0.50001\int_{y_{\min}}^\infty\Big(\frac A{C'y^2}+\frac1{C'^2\delta y^3}\Big)dy\\
&=\ 1.00002\Big(\frac A{C'y_{\min}}+\frac1{2C'^2\delta y_{\min}^2}\Big)\ <\ 1.61\cdot10^{-5}.
\end{align*}

\smallskip\noindent\emph{Kind (iii): many non-rigid columns.}
Let $Y_b:=Y_{iii}\cap H_b$ and $Y_t:=Y_{iii}\cap H_t$. Let $s\in[y_{\min},y_1]$ and $\mathcal W_s:=[(1-\varepsilon)s,s]$, and let $\mathcal Z(s)$
be the set of squares $Z$ with $\Phi(Z)\cap Y_b\cap\mathcal W_s\ne\emptyset$, where $\Phi(Z)=\{y:0<c_Z(y)<1\}$ is as in
Lemma~\ref{lem:A}. For $y\in Y_b\cap\mathcal W_s$ we have $d(y)=y$, and Lemma~\ref{lem:E} gives
$\sum_{Z:\,y\in\Phi(Z)}c_Z(y)\ge(1-\omega_0-E(y))_+$ (the sum is $\ge0$, and $\omega(y)<\omega_0$); every $Z$ in this sum belongs to $\mathcal Z(s)$, meets the line $y\le s$ (so that the height
$v_Z$ of its lowest point satisfies $v_Z\le s$), and has $0<a(Z)<\alpha'(y)\le\alpha'((1-\varepsilon)s)$ (by Lemma~\ref{lem:A}, $\Phi(Z)=\emptyset$ if $a(Z)=0$). Integrating over $Y_b\cap\mathcal W_s$ and
using $\int_{\Phi(Z)}c_Z=\sin a(Z)\cos a(Z)<\alpha'((1-\varepsilon)s)$ (Lemma~\ref{lem:A}), we get
\[
|\mathcal Z(s)|\ \ge\ C'(1-\varepsilon)s\int_{Y_b\cap\mathcal W_s}(1-\omega_0-E(y))_+\,dy .
\]

Let $Z\in\mathcal Z(s)$. As $C'(1-\varepsilon)>C$, we have $a(Z)<\alpha'((1-\varepsilon)s)<\alpha(s)$, so Lemmas~\ref{lem:Q} and~\ref{lem:GZ}
apply to $s$ and $Z$. Every $x\in G_Z^{(s)}$ is either not $s$-rigid from below, or it is $s$-rigid from below and satisfies~(i);
in the latter case Lemma~\ref{lem:Q} would put every ramp point of $Z$ within $w_0$ of an integer, which is impossible because
$Z$ has a ramp point $y\in H$ (as $\Phi(Z)\subset\mathcal R(Z)$). Hence $G_Z^{(s)}\subset\mathcal M'(Z)\cap\mathrm{NR}_s$, and
$|G_Z^{(s)}|>g_0$ by Lemma~\ref{lem:GZ}. Every $Z\in\mathcal Z(s)$ meets a line in $\mathcal W_s$ and has vertical extent
$<1.0001$, and a vertical line through $x$ meets every $Z$ with $x\in G_Z^{(s)}\subset\mathcal M'(Z)$ in a chord of length $\ge1$; these
chords are disjoint and lie in an interval of length $<\varepsilon s+2.01$, so at most $\varepsilon s+2.01$ of the sets $G_Z^{(s)}$ contain
$x$. Hence
\begin{equation}\label{eq:NR}
\mathcal N(s):=|\mathrm{NR}_s|\ \ge\ \frac{g_0\,|\mathcal Z(s)|}{\varepsilon s+2.01}\ \ge\ \frac{g_0C'(1-\varepsilon)s}{\varepsilon s+2.01}
\int_{Y_b\cap\mathcal W_s}(1-\omega_0-E(y))_+\,dy .
\end{equation}
By the reflection $y\mapsto k-y$ and the versions of Lemmas~\ref{lem:Q} and~\ref{lem:GZ} for columns from the top,
$\widetilde{\mathcal N}(s):=|\widetilde{\mathrm{NR}}_s|$ satisfies \eqref{eq:NR} with $Y_b\cap\mathcal W_s$ replaced by $Y_t\cap(k-\mathcal W_s)$.

\smallskip\noindent\emph{Kind (iii): the cost of non-rigid columns.}
For $x\in[0,k]$ let $T^*(x)$ be the least $s\in[y_{\min},y_1]$ with $x\in\mathrm{NR}_s$ (it exists by closedness, if
$x\in\mathrm{NR}_{y_1}$; otherwise $T^*(x)$ is undefined), and define $\widetilde T^*(x)$ in the same way with $\widetilde{\mathrm{NR}}_s$.
By monotonicity, $\{x:T^*(x)\le s\}=\mathrm{NR}_s$ for $s\in[y_{\min},y_1]$; as these sets are finite unions of intervals, $T^*$ is a Borel
function on the Borel set where it is defined, and $\{(x,s):T^*(x)\le s\}$ is a Borel set; likewise for $\widetilde T^*$.
If $T^*(x)=\tau$, then either (W) $\sigma^\tau_x$ has vertical waste $\ge\delta$, or (I) it has vertical waste $<\delta$ and meets a
square of inclination $\ge\alpha(\tau)$; similarly for $\widetilde T^*$ and $\tilde\sigma^\tau_x$. For each $x$ the segments
$\sigma^{T^*(x)}_x\subset\{x\}\times[0,\frac k2-1]$ and $\tilde\sigma^{\widetilde T^*(x)}_x\subset\{x\}\times[\frac k2+1,k]$ are disjoint, so
integrating the vertical waste over $x$ shows that the $x$ of type (W) for $T^*$ and those of type (W) for $\widetilde T^*$
form two sets of total measure $\le W/\delta$; on them $\alpha(T^*),\alpha(\widetilde T^*)\le\alpha(y_{\min})$. Lemma~\ref{lem:Bv}, applied to the
$x$ of type (I) with $\bar y_x=T^*(x)+2$, $\beta(x)=\alpha(T^*(x))$, $\bar y'_x=\widetilde T^*(x)+2$, $\beta'(x)=\alpha(\widetilde T^*(x))$ and
$\bar\beta=\alpha(y_{\min})<2.119\cdot10^{-6}$, bounds the rest. Since $\alpha(\tau)\ge\int_\tau^{y_1}ds/(Cs^2)$, Tonelli's theorem gives
\[
\Big(A+\frac{\alpha(y_{\min})}\delta\Big)W\ \ge\ \int\alpha(T^*(x))\,dx+\int\alpha(\widetilde T^*(x))\,dx\ \ge\ \int_{y_{\min}}^{y_1}\frac{\mathcal N(s)+\widetilde{\mathcal N}(s)}{Cs^2}\,ds .
\]
For $y\in Y_b\subset[y_{\min},(1-\varepsilon)y_1]$ the interval $[y,y/(1-\varepsilon)]$ lies in $[y_{\min},y_1]$, and
\[
\int_y^{y/(1-\varepsilon)}\frac{ds}{s(\varepsilon s+2.01)}\ \ge\ \frac1{1+2.01/(\varepsilon y_{\min})}\int_y^{y/(1-\varepsilon)}\frac{ds}{\varepsilon s^2}
\ =\ \frac1{(1+2.01/(\varepsilon y_{\min}))\,y} .
\]
By \eqref{eq:NR} and Tonelli's theorem again, writing $\int_{Y_b\cap\mathcal W_s}f=\int_{Y_b}f(y)\mathbf 1[y\le s\le y/(1-\varepsilon)]\,dy$, and in
the same way for $\widetilde{\mathcal N}$ (where $d(y)=k-y$),
\[
\Big(A+\frac{\alpha(y_{\min})}\delta\Big)W\ \ge\ \frac{g_0C'(1-\varepsilon)}{C\,(1+2.01/(\varepsilon y_{\min}))}\int_{Y_{iii}}(1-\omega_0-E(y))_+\,\frac{dy}{d(y)} .
\]

\smallskip\noindent\emph{Conclusion.} On $Y_{iii}$ we have $1-\omega_0\le(1-\omega_0-E)_++E\le(1-\omega_0-E)_++\gamma W$, so
$(1-\omega_0)\ell(Y_{iii})\le\int_{Y_{iii}}(1-\omega_0-E)_+\,dy/d(y)+\Gamma W$. Adding the bounds for the three kinds gives
\begin{align*}
(1-\omega_0)\,\ell(H)\ &\le\ \ell(Y_i)+\ell(Y_{ii})+\int_{Y_{iii}}(1-\omega_0-E)_+\frac{dy}{d(y)}+\Gamma W\\
&\le\ \Big(\frac1{\omega_0y_{\min}}+\frac{C'}2A+\frac{C\,(1+2.01/(\varepsilon y_{\min}))}{C'g_0(1-\varepsilon)}\Big(A+\frac{\alpha(y_{\min})}\delta\Big)\\
&\qquad+\Gamma\Big)W .
\end{align*}
With $y_{\min}=236000$ (so $1/(\omega_0y_{\min})<0.021187$, $2.01/(\varepsilon y_{\min})<4.2585\cdot10^{-3}$ and
$\alpha(y_{\min})/\delta<0.21187$), $A<10.6067$, $g_0=0.49997$ and $\Gamma<1.61\cdot10^{-5}$, the bracket is less than
\begin{align*}
&0.021187+\sqrt2\cdot10.6067+\frac{1.0042585}{\sqrt2\cdot0.49997\cdot0.998}\,(10.6067+0.21187)+1.61\cdot10^{-5}\\
&\qquad=\ 30.41797\ldots\ <\ 30.418,
\end{align*}
so, as $(1-\omega_0)h_0\ge0.99977$, $W\ge2\cdot0.99977(\log k-13.06675)/30.418=F(k)$. Thus $W\ge\max(1,F(k))$ for every $k\ge2$. If
$0.0353\log k\le1$, then $W\ge0.0353\log k$ by Corollary~\ref{cor:W1}. Otherwise $\log k>28.32$, and as
$1.99954/30.418-0.0353>0.03043$ and $1.99954\cdot13.06675/30.418<0.85895<0.03043\cdot28.32$, we get
$F(k)-0.0353\log k>0.03043\log k-0.85895>0$. This proves Theorem~\ref{thm:main} with $\kappa=0.0353$. \qed

\section{Remarks}

\begin{remark}[where integrality and closedness enter]
Closedness is used in Lemma~\ref{lem:G} (hence in Lemmas~\ref{lem:A} and~\ref{lem:E} and Corollary~\ref{cor:W1}), in
Propositions~\ref{prop:FL} and~\ref{prop:FLpar} and Lemmas~\ref{lem:slit} and~\ref{lem:crack} (separation of disjoint squares by a line), in
Lemma~\ref{lem:planar} (consecutive squares do not touch), and in
Lemma~\ref{lem:red}, which ties closed packings of $[0,k]^2$ to $s(n)<k$; for squares with disjoint interiors in a
square of side exactly $k$, the $k\times k$ grid shows that nothing can be said. The disjointness of the squares is also used, in the
arguments as given, in Lemmas~\ref{lem:count}, \ref{lem:Kw} and~\ref{lem:Kw9} (centres of disjoint squares are at distance at least $1$;
in Lemma~\ref{lem:count}(a) also the disjointness of $Q_Z$, $S_1$ and $S_2$), in Proposition~\ref{prop:FL}(b) (the regions swept by
the uncovered parts are disjoint parts of the waste), in Lemma~\ref{lem:slit}(a) (the strict inequality $G(y_Z)>1$), in Lemmas~\ref{lem:Bh}
and~\ref{lem:Bv} (the gap regions of different pairs are disjoint parts of the waste), in Lemmas~\ref{lem:Eline} and~\ref{lem:GZ} (chords
of distinct squares at one height are disjoint; in the latter also $q\notin Z$), in Lemma~\ref{lem:Q} and the definition before it
(vertical chords of distinct squares at one abscissa are disjoint), and in Section~6 (disjoint vertical chords bound the multiplicity of
the sets $G_Z^{(s)}$). The integrality of the side enters through
Lemma~\ref{lem:G} (the bound $k-1$ on long chords) and through the columns hanging from the top side: the reflection $y\mapsto k-y$
preserves the heights close to integers only because $k$ is an integer. The quantization of Lemma~\ref{lem:Q} itself only uses one
side, the unit size of the squares and their disjointness.
\end{remark}

\begin{remark}[large inclinations]
The same line argument without the column analysis gives the following dichotomy. Let $k\ge3\cdot10^{10}$ and
$\alpha:=1/(4\sqrt k)\le1.45\cdot10^{-6}$. The lines with $\omega(y)\ge\frac12$ have measure at most $2W$; by Lemma~\ref{lem:Bh} with $\beta\equiv\alpha$, the lines
with $\omega(y)<\frac12$ that meet a square of inclination $\ge\alpha$ have measure at most $AW/(2\alpha)$; and on the remaining lines
Lemma~\ref{lem:A} gives short chords of total length $>1-\omega(y)-0.0313$. Integrating,
\[
0.9687\,k\ \le\ (2.94+1.94\,A\sqrt k)\,W+\sum_{0<a(S)<\alpha}\sin a(S)\cos a(S).
\]
Thus a closed packing with $W=o(\sqrt k)$ must contain at least $(3.87-o(1))\,k^{3/2}$ squares of inclination in $(0,\alpha)$.
This remark is not used in the proofs.
\end{remark}

\begin{remark}[the threshold and the constant as functions of $k$]\label{rem:k}
Put $F(k):=1.99954\,(\log k-13.06675)/30.418$, so that every closed packing satisfies $W=k^2-N\ge\max(1,F(k))$ for every
$k\ge2$ (Section~6), and let $\kappa(k_1):=\inf_{k\ge k_1}\max(1,F(k))/\log k$ be the largest constant that these bounds certify for
all $k\ge k_1$. It is non-decreasing in $k_1$. As the additive constant is negative, $F(k)/\log k=\frac{1.99954}{30.418}(1-13.06675/\log k)$
increases with $k$, towards its limit $1.99954/30.418=0.065735\ldots$, and $F(k)$ exceeds $1$ only for $\log k>28.2792\ldots$; below this
point the certified ratio is $1/\log k$, and beyond it $F(k)/\log k$, which there exceeds $1/28.2793$. Hence
$\kappa(k_1)=0.035361\ldots$ for $2\le k_1\le e^{28.279}$, and $\kappa(k_1)$ increases to $0.065735\ldots$ as $k_1\to\infty$. In contrast
with versions~1.0 and~1.1 of this paper, the constant for all $k$ is thus limited not by the slope but by the point at which $F(k)$ exceeds
$1$, that is, by the additive constant $13.06675\approx\log(2(y_{\min}+1)/(1-\varepsilon))$ coming from the smallest scale $y_{\min}$. A larger
$y_{\min}$ slightly increases the slope but decreases the constant for all $k$ (with unrounded constants, $y_{\min}=3\cdot10^5$ and $10^6$
give slopes $0.0659$ and $0.0663$ and constants $0.0351$ and $0.0338$ for all $k$, all rounded down), while with Lemma~\ref{lem:Bh} as stated $y_{\min}$
cannot be smaller than $1/(C'\cdot1.5\cdot10^{-6})=235702.26\ldots$ (and Lemmas~\ref{lem:Q} and~\ref{lem:GZ} are stated for $s\ge2.36\cdot10^5$;
at $s=235702.26$ the value $w_0=1.265\cdot10^{-5}$ would not suffice). The limit $1/(\sqrt2A_\infty)=\sqrt{Q_*/K_w^*}/(2\sqrt2)\approx0.066667$,
where $A_\infty=2\sqrt{K_w^*/Q_*}$, is the most that this architecture can give for the slope with $K_w^*=9$ and $Q_*=0.32$. For
comparison, versions~1.0 and~1.1 of this paper, which used only the lines with $d(y)\ge\sqrt k/4$, had the bound
$F(k)=(0.99977\log k+0.3819)/30.147$ for $k\ge10^{13}$ and gave $\kappa=0.033$ (its ratio $F(k)/\log k$ decreased to the limit
$0.033163\ldots$), the bound $F(k)=(0.99996\log k+0.386)/36.081$ of an
earlier draft (with $K_w=13$ and the threshold $10^{16}$) gave $\kappa=0.027$, the bound $F(k)=(0.99996\log k+0.386)/37.442$ of an earlier draft
(with $K_w=14$) gave $\kappa=0.026$, the bound
$F(k)=(0.999996\log k+0.386)/51.963$ of an earlier draft (with the multiplicity $K=27$ in place of $K_w$) exceeded $1$ only for
$\log k>51.58$ and gave $\kappa=0.019$, and the bound $F(k)=(0.9987(\log k-1.3884)-0.0313)/84.05$ of an earlier draft exceeded $1$ only
for $\log k>85.58$, and the constant it gave for all $k\ge2$ was $1/85.58\approx0.0117$. Constant lower bounds for $W$ from other sources
raise the constant for all $k$ without changing the slope: with the conditional bound $W\ge4$ for $k\ge6$ of Corollary~\ref{cor:evand}
the infimum becomes
$4/73.91674\ldots=0.054114\ldots$.
\end{remark}

\begin{remark}[the rate]
The proof gives $c^*(k)\gg\log k$, while the upper bounds give $c^*(k)=O(k^{3/5})$. Several heuristics (balancing the
cost of inclinations against the quantization near the walls, and numerical linear-programming bounds for
periodic covering measures) suggest that the truth may be of order $\sqrt k$; we leave this open. The constants in
Theorem~\ref{thm:main} are far from useful for concrete values: since $W$ is an integer, the second bound of the theorem yields
$s(k^2-c)=k$ for every $c\ge1$ as soon as $\log k>15.22\,c+13.07$; for instance (with the exact threshold $13.06675+30.418c/1.99954$) it yields $s(k^2-4)=k$ for $\log k>73.92$ (compare the computer-assisted, unrefereed announcement of $s(k^2-4)=k$ for all
$k\ge5$ in \cite{Dan}), and it says nothing
about $s(12)$ or about any $k$ within computational reach.
\end{remark}

\begin{remark}[the constant]\label{rem:const}
An earlier draft of this argument, which used the lemma of \cite{RV} with its stated constant
$10^{-10}$ and coarser bookkeeping, obtained only $\kappa\approx6\cdot10^{-17}$, and previous drafts of this paper
obtained $\kappa=3\cdot10^{-4}$, $1.9\cdot10^{-3}$, $0.011$, $0.019$, $0.026$, $0.027$ and (versions~1.0 and~1.1) $0.033$. The present constant comes from
Proposition~\ref{prop:FL} (convex gaps, and cuts only where the set of squares met by the chords changes) and its refinement
Proposition~\ref{prop:FLpar} for nearly parallel pairs, which replaces the factor $1/q_0=1/10$ by $Q_*=0.32$; from the angular argument
of Lemma~\ref{lem:Kw} and its computer-assisted refinement Lemma~\ref{lem:Kw9}, which bound the multiplicity of the rectangles at
uncovered points, the only ones that matter, by $K_w=13$ and $K_w^*=9$, the latter sharp for all pairs at distance $<d$ (the
planarity argument of Lemma~\ref{lem:count}(b), with Fodor's theorem \cite{Fod} and Oler's inequality \cite{Ole,FG}, gives $K=27$ at every
point and is no longer needed); from Lemma~\ref{lem:slit}(a), by which each pair is charged to a single slit;
from capping the inclinations and splitting off the lines with a wide gap in Lemma~\ref{lem:Bh}; from Lemmas~\ref{lem:A} and~\ref{lem:E}, by which the
short chords on a line with little waste add up to almost a full unit (this halves the cost of the columns); from
Lemma~\ref{lem:Eline}, by which the excess of the long chords is paid for by the waste, so that the lines down to the constant
distance $y_{\min}$ from the walls can be used (this doubles the logarithmic range of scales once more, and with it the slope); from
Lemma~\ref{lem:crack}, by which a rigid column through the end of a square passes through a tight crack, so that the exceptional
columns under a square have measure about $1/C$ rather than $2/C$; from measuring heights from both the bottom and the top side,
which doubles the logarithmic range of scales; and from working with a continuum of scales instead of dyadic bands. With
$g_0\approx1-\frac1C$ the bracket in the proof is about $(\frac{C'}2+\frac C{C'(1-1/C)})A$, which is minimal, equal to $2\sqrt2\,A$, at
$C=2$, $C'=2\sqrt2$. The factor $\sqrt{K_w^*/Q_*}\approx5.3$ in $A$ comes from using Propositions~\ref{prop:FL}
and~\ref{prop:FLpar} for pairs that are consecutive only on a short set of lines, typically at corners where several squares meet,
whereas in simple configurations, such as a single slightly tilted square in a grid, the ratio of $\int_Y2\beta$ to $W$ in
Lemma~\ref{lem:Bh} is of order $1$. A lemma charging such short-lived pairs to the slits of neighbouring long-lived pairs with a
bounded constant would improve $\kappa$ further; we have not been able to prove one. In numerical experiments (randomized dense packings and targeted searches) the largest multiplicity of the rectangles
$R_e$ found by the accompanying search programs was $5$ (against $27$; an earlier search, whose program was not preserved,
reported $8$); and the largest number of pairs of squares at distance $<10^{-4}$ (whether or not they
are edges of $G_h$ or $G_v$) whose rectangles $R(P,Q)$ contain a common uncovered point was $9$, equal to $K_w^*$, in a degenerate
configuration: two axis-parallel squares on both sides of a slit of width about $10^{-20}$ through the point, each nearly touching
two slightly tilted squares on either side (random searches found at most $6$). We also found no configuration with $\sum\lambda_r^2<\frac13$ in Proposition~\ref{prop:FLpar}
(against $0.32$; the smallest value found was $0.333335$). The search programs and their results are part of the accompanying material. With $Q_*=\frac13$ the same proof (with Lemma~\ref{lem:Bh} stated for $\bar\beta\le1.4999\cdot10^{-6}$ and Lemma~\ref{lem:Bv} for
$\bar\beta\le2.99\cdot10^{-6}$, which still cover their uses in Section~6) would give
the slope $0.0670$ and the constant $0.0357$ for all $k$ (both rounded down).
Only the numerical constants depend on the computer-assisted Lemma~\ref{lem:Kw9}: with the analytic bound $K_w=13$ of
Lemma~\ref{lem:Kw} in place of $K_w^*$, the proof of Section~6 gives $A<12.7476$, a bracket less than $36.493$ (with
$\Gamma<1.93\cdot10^{-5}$), the slope $1.99954/36.493>0.0547$ and $\kappa=0.0319$. Even with $A$ of order $1$, this
architecture would limit the slope to roughly $\frac1{\sqrt2A}$, and, with the present $y_{\min}$, the constant for all $k$ to less than
$1/13.06$, so that a bound useful for concrete $k$ would require a different rate rather than better constants.
\end{remark}

\subsection*{Changes from version 1.1}
The lines near the walls are now used. Lemma~\ref{lem:E} (short chords against the excess $E(y)$ of the long chords) and
Lemma~\ref{lem:Eline} (a bound of $E(y)$ by the waste on lines whose squares have small inclination) are new; in Section~6 the
cutoff $d(y)\ge\sqrt k/4$ is replaced by the constant $d(y)\ge y_{\min}=236000$, the term $\eta(y)$ is replaced by $E(y)$, and the
resulting integral $\Gamma$ enters the bracket; the second bound of the theorem now holds for every $k\ge2$ (for $k<10^{12}$ it follows
from Corollary~\ref{cor:W1}), so the threshold $k_2=10^{13}$ is no longer needed. Lemma~\ref{lem:Bv} is stated for
$\bar\beta\le3\cdot10^{-6}$ (its proof is unchanged), and Lemmas~\ref{lem:Q} and~\ref{lem:GZ} for $s\ge2.36\cdot10^5$, with
$w_0=1.265\cdot10^{-5}$ and $g_0=0.49997$. The constant improves from $\kappa=0.033$ to $\kappa=0.0353$, and the asymptotic slope from
$0.0331$ to $0.0657$; the numbers derived from them (the threshold for $s(k^2-4)=k$, the constant without Lemma~\ref{lem:Kw9}, and
Remarks~\ref{rem:k} and~\ref{rem:const}) are updated accordingly. Corollary~\ref{cor:evand} is new: it records the constant
$0.0541$ for all $k\ge2$ that follows if one also uses the unrefereed computer-assisted result $s(k^2-3)=k$ ($k\ge6$) of \cite{Dan};
the references \cite{Lev,Wan} are new, and the entry \cite{Dan} lists the further files used. Following the findings of an AI review of
version~1.1 by the Squares Project \cite{LevR}, the following passages were also revised, without changing the results: the standing
assumption of Section~4 and the proof of Lemma~\ref{lem:Kw9} now agree that $k\ge4$ there, and Section~4 states that Sections~4 and~5 are
used in the proof of Theorem~\ref{thm:main} only for $k\ge10^{12}$; the proof of Lemma~\ref{lem:Kw9} now explains why the extension of the
angle intervals to $2\pi$ covers the arc $(\widetilde{2\pi},2\pi)$, shows that the floating-point count of points is exact, says which
parts of the search programs the re-verification uses, describes the authorship of the two further coverage checks, and reports the
replay of the computation by the Squares Project; the first remark of Section~7 lists further uses of disjointness; and the remark after
Lemma~\ref{lem:Kw9} and the remark on large inclinations state that they are not used in the proofs. The new reference \cite{LevR} records
that review. Nothing else in the mathematics has changed.

\subsection*{Acknowledgements}
I thank J.~Levy and the Squares Project for an AI review of version~1.1 of this paper and for replaying the computation of
Lemma~\ref{lem:Kw9} \cite{LevR}; this version takes their findings into account.

\subsection*{Use of AI}
Claude (Anthropic) was used extensively in developing the arguments, writing the text and the programs (including the
computation for Lemma~\ref{lem:Kw9}), and in checking them. It is not an author; the author takes full responsibility for the
contents. The paper has not been refereed.

\subsection*{Licence}
This text is distributed under the Creative Commons Attribution 4.0 International licence (CC BY 4.0); the accompanying programs and data are
distributed under the MIT licence.

\begin{thebibliography}{99}
\raggedright
\bibitem{AB} M.~Z.~Arslanov, H.~D.~Bui, Note on ``Efficient packings of unit squares in a large square'',
\emph{Discrete Comput.\ Geom.} \textbf{76} (2026), no.~2, 805--812 (published online 3 August 2025), doi:10.1007/s00454-025-00767-w.
\bibitem{AMS} M.~Z.~Arslanov, S.~A.~Mustafin, Z.~K.~Shangitbayev, Improved packings of $n(n-1)$ unit squares in a square,
\emph{Electron.\ J.\ Combin.} \textbf{28}(4) (2021), \#P4.22, doi:10.37236/8586.
\bibitem{Ben10} W.~Bentz, Optimal packings of 13 and 46 unit squares in a square, \emph{Electron.\ J.\ Combin.} \textbf{17}
(2010), \#R126, doi:10.37236/398.
\bibitem{Ben16} W.~Bentz, Optimal packings of 22 and 33 unit squares in a square, arXiv:1606.03746 (2016).
\bibitem{BuiGood} H.~D.~Bui, Square packing with asymptotically smallest waste only needs good squares,
arXiv:2504.09489 (2025).
\bibitem{Bui} H.~D.~Bui, Square packing with $O(x^{0.6})$ wasted area, arXiv:2508.04603v2 (2025, revised 2026).
\bibitem{Che} chelokot, \emph{square-packing-archive}, GitHub repository,
\url{https://github.com/chelokot/square-packing-archive}, files \texttt{docs/nagamochi-score-counterexample.md} and
\texttt{docs/nagamochi-compensation-proof.md}, commit \texttt{753079e} (accessed 3 October 2026).
\bibitem{CG09} F.~Chung, R.~Graham, Packing equal squares into a large square, \emph{J.\ Combin.\ Theory Ser.\ A}
\textbf{116} (2009), no.~6, 1167--1175, doi:10.1016/j.jcta.2009.02.005.
\bibitem{CG} F.~Chung, R.~Graham, Efficient packings of unit squares in a large square, \emph{Discrete Comput.\ Geom.}
\textbf{64} (2020), no.~3, 690--699, doi:10.1007/s00454-019-00088-9.
\bibitem{Dan} E.~Daniel, \emph{square-packing}, GitHub repository, \url{https://github.com/evand/square-packing},
directories \texttt{s12/certificates/k2m3} and \texttt{s12/certificates/k2m4}, files \texttt{s12/search/FRIEDMAN.md},
\texttt{s12/CREDITS.md} and \texttt{s12/docs/k2m3.html}, and Lean files \texttt{s12/lean/Sqpack/Bentz.lean} and
\texttt{s12/lean/Sqpack/ValidSplit7.lean}, commit \texttt{7ff3b21} (accessed 5 October 2026; these files are unchanged at commit
\texttt{13ee36e}, accessed 6 October 2026).
\bibitem{EG} P.~Erd\H{o}s, R.~L.~Graham, On packing squares with equal squares, \emph{J.\ Combin.\ Theory Ser.\ A}
\textbf{19} (1975), no.~1, 119--123, doi:10.1016/0097-3165(75)90099-0.
\bibitem{FLINT} The FLINT team, \emph{FLINT: Fast Library for Number Theory}, version 3.6.0 (2026), \url{https://flintlib.org}.

\bibitem{Fod} F.~Fodor, The densest packing of 12 congruent circles in a circle, \emph{Beitr\"age Algebra Geom.}
\textbf{41} (2000), no.~2, 401--409.
\bibitem{FG} J.~H.~Folkman, R.~L.~Graham, A packing inequality for compact convex subsets of the plane,
\emph{Canad.\ Math.\ Bull.} \textbf{12} (1969), no.~6, 745--752, doi:10.4153/CMB-1969-096-7.
\bibitem{Fri} E.~Friedman, Packing unit squares in squares: a survey and new results, \emph{Electron.\ J.\ Combin.},
Dynamic Survey DS7 (1998, revised 2009), doi:10.37236/28.
\bibitem{Arb} F.~Johansson, Arb: efficient arbitrary-precision midpoint-radius interval arithmetic, \emph{IEEE Trans.\ Comput.}
\textbf{66} (2017), no.~8, 1281--1292, doi:10.1109/TC.2017.2690633.

\bibitem{pyflint} F.~Johansson, O.~Benjamin, et al., \emph{python-flint}, version 0.9.0 (2026), \url{https://github.com/flintlib/python-flint}.

\bibitem{Kar} H.~Karaku\c{s}, A counterexample to Nagamochi's scoring lemma and a new rectangle packing bound,
arXiv:2609.37410 (2026).
\bibitem{KS} M.~J.~Kearney, P.~Shiu, Efficient packing of unit squares in a square, \emph{Electron.\ J.\ Combin.} \textbf{9}
(2002), \#R14, doi:10.37236/1631.
\bibitem{Lev} J.~Levy, \emph{The Squares Project}, GitHub repository, \url{https://github.com/jlevy/squares}, files
\texttt{packing/frontier/RESULTS.md}, \texttt{packing/frontier/results.yaml} (entry T-064) and \texttt{packing/frontier/evidence.yaml},
commit \texttt{ebf2327} (accessed 6 October 2026).
\bibitem{LevR} J.~Levy, \emph{The Squares Project}, GitHub repository, \url{https://github.com/jlevy/squares}: AI review of version~1.1 of
this paper, file \texttt{docs/project/reviews/review-2026-10-05-squarepacker-k2-minus-c.md}, commit \texttt{6953107}, and issue~\#368
(comments of 5 and 6 October 2026, UTC; accessed 7 October 2026).
\bibitem{McC} R.~McClenagan, Optimally packing a large square by unit squares, arXiv:2602.01484 (2026).
\bibitem{Nag} H.~Nagamochi, Packing unit squares in a rectangle, \emph{Electron.\ J.\ Combin.} \textbf{12} (2005), \#R37, doi:10.37236/1934.
\bibitem{Ole} N.~Oler, An inequality in the geometry of numbers, \emph{Acta Math.} \textbf{105} (1961), 19--48,
doi:10.1007/BF02559533.
\bibitem{RV} K.~F.~Roth, R.~C.~Vaughan, Inefficiency in packing squares with unit squares, \emph{J.\ Combin.\ Theory Ser.\ A}
\textbf{24} (1978), no.~2, 170--186, doi:10.1016/0097-3165(78)90005-5.
\bibitem{Wan} wand125, \emph{valid7-independent-check}, GitHub repository,
\url{https://github.com/wand125/valid7-independent-check}, commit \texttt{c561dbb}, with run records in release \texttt{records-v1}
(accessed 6 October 2026).
\end{thebibliography}

\end{document}
